%% Final report of the MO422 project: a from-scratch C++ implementation of TFHE.
%%
%% Template: SBC (sbc-template.sty, sbc.bst and caption2.sty, distributed by the
%% Brazilian Computer Society and included in this folder without changes).
%%
%% Build:  make        (pdflatex; bibtex; pdflatex x2)
%% Data:   data/results.tex and data/*.dat are produced by experiments/summarize.py
%%         from the raw measurements in data/raw/ (see experiments/run_all.sh).
\documentclass[12pt]{article}

\usepackage{sbc-template}
\usepackage[T1]{fontenc}
\usepackage[utf8]{inputenc}
\usepackage[brazilian,english]{babel}
\usepackage{graphicx,url}
\usepackage{amsmath,amssymb,amsthm}
\usepackage{booktabs}
\usepackage{xcolor}
\usepackage{tikz}
\usepackage{pgfplots}
\usepackage{algorithm}
\usepackage{algpseudocode}
\usepackage{listings}
\usepackage{pifont}
\usepackage[hidelinks]{hyperref}
\usetikzlibrary{positioning,decorations.pathreplacing,calc,arrows.meta}
\usepgfplotslibrary{groupplots}
% Overlay specifications in the shared TikZ figures only matter in the slides.
\tikzset{visible on/.style={}}
\pgfplotsset{compat=1.18}

\definecolor{cBlue}{HTML}{2C5F8A}
\definecolor{cOrange}{HTML}{C9762B}
\definecolor{cGreen}{HTML}{3B8B5C}
\definecolor{cRed}{HTML}{B03A3A}
\definecolor{cGray}{HTML}{6E6E6E}
\definecolor{cNavy}{HTML}{1E2B3C}

\lstset{basicstyle=\ttfamily\fontsize{9}{11}\selectfont, columns=fullflexible, keepspaces=true,
        language=C++, keywordstyle=\color{cBlue}\bfseries,
        commentstyle=\color{cGray}\itshape, frame=lines, framesep=4pt,
        rulecolor=\color{cGray!60}, xleftmargin=0.3cm, xrightmargin=0.3cm,
        aboveskip=8pt, belowskip=4pt, numbers=none}

\theoremstyle{plain}
\newtheorem{proposition}{Proposition}
\newtheorem{lemma}{Lemma}
\theoremstyle{definition}
\newtheorem{definition}{Definition}
\newtheorem{example}{Example}

\sloppy
\raggedbottom   % keep paragraph spacing uniform when a figure floats to the next page
\setlength{\belowcaptionskip}{4pt}

% Generated by experiments/summarize.py -- do not edit.
\newcommand{\OurNandMean}{22.9}
\newcommand{\SlowdownSpq}{2.3}
\newcommand{\SpeedupNay}{4.9}
\newcommand{\TotalGates}{12{,}500}
\newcommand{\OurBenchGates}{2{,}500}
\newcommand{\OrigGates}{3{,}000}
\newcommand{\OrigGatesEach}{1{,}500}
\newcommand{\OurRuns}{5}
\newcommand{\OrigRuns}{3}
\newcommand{\OurKS}{0.8}
\newcommand{\OurNandRange}{22.1--23.8}
\newcommand{\SpqNandRange}{9.7--9.8}
\newcommand{\NayNandRange}{107.8--122.1}
\newcommand{\OurKeygenRange}{0.21--0.22}
\newcommand{\SpqKeygenRange}{0.53}
\newcommand{\NayKeygenRange}{0.88--0.94}
\newcommand{\OurVarRange}{0.78--0.92}
\newcommand{\SpqVarRange}{1.00--1.12}
\newcommand{\NayVarRange}{1.00--1.16}
\newcommand{\OurMeanMax}{\ensuremath{1.2\cdot10^{-3}}}
\newcommand{\SpqMeanMax}{\ensuremath{2.0\cdot10^{-4}}}
\newcommand{\NayMeanMax}{\ensuremath{1.9\cdot10^{-4}}}
\newcommand{\OurRSS}{143}
\newcommand{\SpqRSS}{253}
\newcommand{\NayRSS}{253}
\newcommand{\BkMB}{59.1}
\newcommand{\KskMB}{78.9}
\newcommand{\AdderEight}{0.72}
\newcommand{\AdderSixteen}{1.43}
\newcommand{\AdderTwentyFour}{2.16}
\newcommand{\AdderEightText}{$173+94=267$ (8 bits)}
\newcommand{\AdderSixteenText}{$40000+30000=70000$ (16 bits)}
\newcommand{\AdderTwentyFourText}{$9999999+7777777=17777776$ (24 bits)}
\newcommand{\FftNaive}{325}
\newcommand{\FftFast}{5.4}
\newcommand{\FftSpeedup}{60}
\newcommand{\BootSdIn}{\ensuremath{1.0\cdot10^{-2}}}
\newcommand{\BootSdOut}{\ensuremath{2.9\cdot10^{-3}}}
\newcommand{\BootSdInRaw}{1.030000e-02}
\newcommand{\BootSdOutRaw}{2.850000e-03}
\newcommand{\BRShare}{96.5}
\newcommand{\FftShare}{42}
\newcommand{\MacShare}{22}
\newcommand{\FreqShare}{63}
\newcommand{\DecShare}{7}
\newcommand{\OtherShare}{25}
\newcommand{\ExtHot}{21.8}
\newcommand{\ExtKey}{33.6}
\newcommand{\ExtParts}{18.6}
\newcommand{\ExtAlloc}{3}
\newcommand{\MacHot}{0.11}
\newcommand{\MacKey}{0.64}
\newcommand{\MacRatio}{5.6}
\newcommand{\MemShare}{34}
\newcommand{\BoundBR}{\ensuremath{1.4\cdot10^{-5}}}
\newcommand{\BoundTotal}{\ensuremath{2.2\cdot10^{-5}}}
\newcommand{\AvgBR}{\ensuremath{4.7\cdot10^{-6}}}
\newcommand{\AvgKS}{\ensuremath{4.3\cdot10^{-6}}}
\newcommand{\AvgTotal}{\ensuremath{9.0\cdot10^{-6}}}
\newcommand{\BiasStd}{\ensuremath{1.2\cdot10^{-3}}}
\newcommand{\NoiseKeys}{10}
\newcommand{\NoiseGates}{1{,}000}
\newcommand{\NoiseVarBR}{\ensuremath{4.6\cdot10^{-6}}}
\newcommand{\NoiseVarKS}{\ensuremath{9.0\cdot10^{-6}}}
\newcommand{\NoiseBiasStd}{\ensuremath{9.8\cdot10^{-4}}}
\newcommand{\CPU}{Intel Core Ultra 7 256V}
\newcommand{\Compiler}{GCC 15.3.1}
\newcommand{\LibLines}{526}
\newcommand{\AvgTotalRaw}{8.985973e-06}
\newcommand{\NoiseVarKSRaw}{9.036360e-06}
\newcommand{\BiasStdRaw}{1.196040e-03}


\newcommand{\T}{\mathbb{T}}
\newcommand{\Z}{\mathbb{Z}}
\newcommand{\R}{\mathbb{R}}
\newcommand{\B}{\mathbb{B}}
\newcommand{\TN}{\T_N[X]}
\newcommand{\ZN}{\Z_N[X]}
\newcommand{\BN}{\B_N[X]}
\newcommand{\round}[1]{\left\lfloor #1 \right\rceil}
\newcommand{\Bg}{B_g}
\newcommand{\Bks}{B_{\mathrm{ks}}}
\newcommand{\BK}{\mathrm{BK}}
\newcommand{\KSK}{\mathrm{KSK}}
\newcommand{\ext}{\boxdot}
\newcommand{\Dec}{\mathrm{Dec}}
\newcommand{\code}[1]{\texttt{#1}}

\title{TFHE from Scratch: A C++ Implementation of\\Gate Bootstrapping over the Torus}

\author{Matheus de Oliveira Saldanha}

\address{Instituto de Computação -- Universidade Estadual de Campinas (Unicamp)\\
  Campinas -- SP -- Brazil
  \email{matheus.saldanha@dac.unicamp.br}}

\begin{document}

\maketitle

\begin{abstract}
Fully homomorphic encryption (FHE) lets a server compute on data that it
cannot read. TFHE does this one Boolean gate at a time: after every gate, a
procedure called bootstrapping cleans the noise of the result, so circuits of
any depth can be evaluated. In this project we implemented TFHE gate
bootstrapping from scratch in C++, using the same 128-bit parameters as the
original TFHE library. This report explains the scheme step by step, from LWE
encryption to a working encrypted adder, describes the main implementation
choices (the negacyclic FFT and a trick that computes signed digits with a
single addition), and compares our code with the original library. A NAND gate
takes \OurNandMean\,ms in our implementation, which is $\SlowdownSpq\times$
slower than the original library with its assembly FFT and $\SpeedupNay\times$
faster than its portable version. We saw no decryption errors in \TotalGates{}
gates, and the measured noise matches an average-case analysis.
\end{abstract}

\begin{resumo}
\selectlanguage{brazilian}%
A criptografia totalmente homomórfica (FHE) permite que um servidor faça
cálculos sobre dados que ele não consegue ler. O TFHE faz isso uma porta lógica
de cada vez: depois de cada porta, um procedimento chamado \emph{bootstrapping}
limpa o ruído do resultado, e assim circuitos de qualquer profundidade podem
ser avaliados. Neste projeto implementamos o \emph{gate bootstrapping} do TFHE
do zero em C++, com os mesmos parâmetros de 128 bits da biblioteca original.
Este relatório explica o esquema passo a passo, da cifra LWE até um somador
cifrado funcionando, descreve as principais escolhas de implementação (a FFT
negacíclica e um truque que calcula dígitos com sinal com uma única soma) e
compara o nosso código com a biblioteca original. Uma porta NAND leva
\OurNandMean\,ms na nossa implementação, $\SlowdownSpq$ vezes mais que a
biblioteca original com a FFT em assembly e $\SpeedupNay$ vezes menos que a sua
versão portável. Não observamos nenhum erro de decifração em \TotalGates{}
portas, e o ruído medido concorda com uma análise de caso médio.
\selectlanguage{english}%
\end{resumo}

%% ===========================================================================
\section{Introduction}
\label{sec:intro}

The idea of computing on encrypted data was proposed by
\cite{rivest1978databanks}, but the first scheme that supports any computation
only came thirty years later, with the work of \cite{gentry2009fhe}. Today,
practical fully homomorphic encryption (FHE) schemes are based on lattice
problems, mainly Learning With Errors (LWE) \cite{regev2005lwe} and its
version with polynomials \cite{lyubashevsky2010ringlwe}. In these schemes every
ciphertext contains a small random error, called noise, which is needed for
security. Each operation on ciphertexts makes the noise larger, and when it
gets too large the message can no longer be decrypted. Gentry's solution was
\emph{bootstrapping}: an operation that takes a noisy ciphertext and produces a
new ciphertext of the same message with small noise.

The first bootstrapping procedures took minutes. Building on the GSW scheme
\cite{gentry2013gsw}, the FHEW scheme \cite{ducas2015fhew} reduced this to
less than a second by bootstrapping after every Boolean gate. TFHE, proposed by
\cite{chillotti2016faster} and described in full by \cite{chillotti2020tfhe},
improved this approach: it works over the real torus $\T=\R/\Z$ and introduces an \emph{external product}
between two kinds of ciphertexts. With it, a bootstrapped gate takes tens of
milliseconds on one processor core, and the public key that the server needs
shrinks from gigabytes to tens of megabytes. Because every gate gives back a
ciphertext with fresh, small noise, any circuit can be evaluated, no matter how
deep.

The goal of this project was to study TFHE, implement it from scratch in C++,
and compare our implementation with the original library \cite{tfhelib} in
four aspects: correctness, performance, memory and noise. Our
implementation\footnote{The source code, the tests and the scripts that produce
every number in this report are in the project repository
(\texttt{github.com/zinho02/learning-tfhe}).} has no external dependencies and uses exactly the
same parameters as the original library. In this report we:
\begin{itemize}
  \item explain the scheme step by step, showing for each part which problem
        it solves, with small examples (Sections~\ref{sec:background}
        and~\ref{sec:bootstrapping});
  \item describe the implementation choices that matter most for correctness
        and speed (Section~\ref{sec:impl});
  \item compare our implementation with the original library and check that
        the measured noise agrees with the theory (Section~\ref{sec:eval}).
\end{itemize}
Section~\ref{sec:conclusion} concludes.

\paragraph{How the scheme is presented.}
TFHE uses several kinds of ciphertexts and operations, and each one exists to
solve a problem created by the others. Instead of defining all of them at the
start, Sections~\ref{sec:background} and~\ref{sec:bootstrapping} introduce them
in the order in which the problems appear (Figure~\ref{fig:roadmap}): first an
encryption scheme with noise (LWE), then its two limitations, then the plan
TFHE uses to fix both (looking up a value in a table by rotating it), then the
two tools that this plan needs (TRLWE and the external product), and finally
the five steps of a bootstrapped gate. Readers interested only in the
implementation and in the results can go directly to Section~\ref{sec:impl}.

% Roadmap of the concepts in Sections 2 and 3.
\begin{figure}[!htb]
\centering
% The chain of concepts followed by the report (Sections 2 and 3).
\begin{tikzpicture}[font=\footnotesize, >=Stealth,
    con/.style={draw=cNavy!45, rounded corners=3pt, align=center, inner sep=3pt,
                minimum height=1.05cm, text width=3.25cm},
    stp/.style={draw=cNavy!45, rounded corners=3pt, align=center, inner sep=3pt,
                minimum height=0.95cm, text width=2.35cm, fill=cOrange!10, font=\scriptsize},
    sec/.style={font=\scriptsize\bfseries, cGray, roadmap sec/.try},
    flow/.style={->, line width=0.7pt, cNavy!70}]
  \node[con, fill=cGray!8]  (torus) at (0,0)    {\textbf{Torus}\\points on a circle;\\noise removed by\\rounding};
  \node[con, fill=cBlue!8]  (lwe)   at (3.7,0)  {\textbf{LWE}\\mask + noise;\\supports addition};
  \node[con, fill=cRed!7]   (lim)   at (7.4,0)  {\textbf{Two problems}\\no products,\\noise keeps growing};
  \node[con, fill=cGreen!9] (plan)  at (11.1,0) {\textbf{The plan}\\bootstrapping as a\\rotated lookup table};
  \node[con, fill=cBlue!8]  (trlwe) at (7.4,-1.75)  {\textbf{TRLWE}\\tables that can\\be rotated};
  \node[con, fill=cBlue!8]  (trgsw) at (11.1,-1.75) {\textbf{TRGSW, CMux}\\multiply by an\\encrypted bit};
  \node[stp] (ms) at (0,-3.6)    {modulus switching,\\test polynomial};
  \node[stp] (br) at (2.85,-3.6) {blind rotation};
  \node[stp] (se) at (5.7,-3.6)  {sample extraction};
  \node[stp] (ks) at (8.55,-3.6) {key switching};
  \node[stp, fill=cGreen!12] (gt) at (11.4,-3.6) {gates and\\an adder};
  \draw[flow] (torus) -- (lwe);
  \draw[flow] (lwe) -- (lim);
  \draw[flow] (lim) -- (plan);
  \draw[flow] (plan.south) ++(-0.9,0) -- ++(0,-0.2) -| (trlwe.north);
  \draw[flow] (plan) -- (trgsw);
  \draw[flow] (trlwe.west) -| (br.north);
  \draw[flow] (trgsw.south) |- ++(0,-0.35) -| ([xshift=4pt]br.north);
  \draw[flow] (ms) -- (br);
  \draw[flow] (br) -- (se);
  \draw[flow] (se) -- (ks);
  \draw[flow] (ks) -- (gt);
  \node[sec, anchor=south west] at (torus.north west) {\S2.1};
  \node[sec, anchor=south west] at (lwe.north west) {\S2.2};
  \node[sec, anchor=south west] at (lim.north west) {\S2.3};
  \node[sec, anchor=south west] at (plan.north west) {\S2.4};
  \node[sec, anchor=south west] at (trlwe.north west) {\S2.5};
  \node[sec, anchor=south east] at (trgsw.north east) {\S2.6};
  \node[sec, anchor=north] at (5.7,-4.15) {\S3: one gate, step by step};
\end{tikzpicture}

\caption{How the scheme is presented. Each concept of
Section~\ref{sec:background} answers a question left open by the one before
it. The two tools of Sections~\ref{sec:trlwe} and~\ref{sec:extprod} are what
the plan of Section~\ref{sec:plan} needs, and Section~\ref{sec:bootstrapping}
puts everything together in the five steps of a bootstrapped gate.}
\label{fig:roadmap}
\end{figure}


%% ===========================================================================
\section{Background: From Noisy Encryption to an Encrypted Table Lookup}
\label{sec:background}

This section builds, one step at a time, the tools that TFHE combines into a
bootstrapped gate. Each subsection answers a question left open by the one
before it. We follow the presentation of \cite{chillotti2020tfhe}, in the
setting used by the original library; the guide by \cite{joye2022guide}
explains the same material in more detail. We write $\B=\{0,1\}$, $\round{x}$
for the integer closest to $x$, and $\vartheta=\sigma^2$ for the variance of a
random variable.

\subsection{The Torus}
\label{sec:torus}

All ciphertexts in this report work the same way: whoever has the key can
compute the message plus a small error, and then removes the error by
rounding. For this we need a set where messages can be added, where a small
error can be removed by rounding to the nearest message, and where additions
never overflow. The real torus $\T=\R/\Z$, the real numbers modulo $1$, has all
three properties. It is easiest to think of it as a circle of length one
(Figure~\ref{fig:torus}): adding a number moves us along the circle, and after
a full turn we are back at the same point, so $0.25$, $1.25$ and $-0.75$ are the
same element of $\T$.

\paragraph{Messages and rounding.}
A message $m\in\Z_M$ is placed at the point $m/M$, so the $M$ possible messages
are evenly spread around the circle. To decode a noisy point, we round it to
the closest of these $M$ points. This gives the right answer whenever the error
is smaller than half the distance between two messages, that is, smaller than
$1/(2M)$. For example, with $M=8$ every point closer than $1/16$ to $1/8$ is
decoded as $1/8$ (Figure~\ref{fig:torus}a). The logic gates of
Section~\ref{sec:gates} use an even simpler encoding: \textsf{true} is
$+\frac18$, \textsf{false} is $-\frac18$, and to decode we only check in which
half of the circle the point is, which allows errors up to $\frac18$
(Figure~\ref{fig:torus}b).

% The torus as a circle: (a) messages of Z_8 and the decoding region of 1/8;
% (b) bit encoding +-1/8 and the sign-based decision used for gates.
\begin{figure}[!htb]
\centering
\begin{tikzpicture}[font=\small]
  \def\R{1.75}
  % ---------------- (a) ----------------
  \begin{scope}
    \fill[cBlue!22] (0,0) -- (22.5:\R) arc (22.5:67.5:\R) -- cycle;
    \draw[line width=0.9pt] (0,0) circle (\R);
    \foreach \k/\lab in {0/0, 1/\tfrac18, 2/\tfrac14, 3/\tfrac38, 4/\tfrac12,
                         5/\tfrac58, 6/\tfrac34, 7/\tfrac78} {
      \fill (\k*45:\R) circle (1.8pt);
      \node at (\k*45:\R+0.42) {$\lab$};
    }
    \draw[cBlue, densely dashed] (0,0) -- (22.5:\R) (0,0) -- (67.5:\R);
    \node[font=\footnotesize, align=center, cBlue!80!black] at (45:0.95)
      {decodes\\to $\tfrac18$};
    \node at (-\R-0.6, \R+0.35) {(a)};
  \end{scope}
  % ---------------- (b) ----------------
  \begin{scope}[xshift=7.6cm]
    \fill[cBlue!14] (0,0) -- (\R,0) arc (0:180:\R) -- cycle;
    \fill[cOrange!16] (0,0) -- (-\R,0) arc (180:360:\R) -- cycle;
    \draw[line width=0.9pt] (0,0) circle (\R);
    \draw[cGray, densely dotted] (-\R-0.15,0) -- (\R+0.15,0);
    \node[right, font=\footnotesize] at (\R+0.15,0) {$0$};
    \node[left, font=\footnotesize] at (-\R-0.15,0) {$\tfrac12$};
    \fill[cBlue] (45:\R) circle (2.4pt);
    \fill[cOrange] (-45:\R) circle (2.4pt);
    \node[cBlue!80!black, anchor=south west] at (45:\R+0.05) {$+\tfrac18$: \textsf{true}};
    \node[cOrange!80!black, anchor=north west] at (-45:\R+0.05) {$-\tfrac18$: \textsf{false}};
    \node[font=\footnotesize, align=center] at (0,0.75) {phase in $(0,\tfrac12)$\\$\Rightarrow$ \textsf{true}};
    \node[font=\footnotesize, align=center] at (0,-0.75) {phase in $(-\tfrac12,0)$\\$\Rightarrow$ \textsf{false}};
    \node at (-\R-0.6, \R+0.35) {(b)};
  \end{scope}
\end{tikzpicture}
\caption{The torus $\mathbb{T}=\mathbb{R}/\mathbb{Z}$ drawn as a circle of
length one. (a) The eight messages of $\mathbb{Z}_8$, placed at the multiples
of $1/8$; the shaded slice contains the phases that decode to $1/8$, that is,
all points at distance less than $1/16$. (b) The bit encoding used by the
gates: \textsf{true} is $+1/8$ and \textsf{false} is $-1/8$. Decryption only
checks in which half of the circle the phase is, so it allows an error of up
to $1/8$ on either side.}
\label{fig:torus}
\end{figure}


\paragraph{Integers times torus elements, but not torus times torus.}
We can multiply a point $t$ of the torus by an integer $k$: $k\cdot t$ is just
$t$ added $k$ times. But we cannot multiply two points of the torus, because the
result would depend on which real number we use to represent each point. For
example, $\frac12$ and $\frac32$ are the same point, but
$\frac12\cdot\frac12=\frac14$ while $\frac32\cdot\frac12=\frac34$. This limitation
has a big influence on the whole scheme: TFHE never multiplies two ciphertexts,
and Section~\ref{sec:extprod} explains how it still manages to multiply a
ciphertext by an encrypted bit.

\subsection{LWE Encryption}
\label{sec:lwe}

\paragraph{A first attempt.}
Let the secret key be a list of $n$ random bits, $s\in\B^n$. To hide a point
$\mu\in\T$, we choose a random public vector $\mathbf a\in\T^n$ and publish the
pair $(\mathbf a,b)$ with $b=\langle\mathbf a,s\rangle+\mu$. Only the owner of
$s$ can compute the \emph{mask} $\langle\mathbf a,s\rangle=\sum_ia_is_i$ and
subtract it from $b$ to recover $\mu$. This scheme, however, is easy to break.
Every ciphertext of a known message (for example, of $0$) gives the attacker a
linear equation in the $n$ unknown bits of $s$, and with $n$ such equations the
attacker finds $s$ by Gaussian elimination.

\paragraph{Adding noise.}
\cite{regev2005lwe} fixed this by making every equation slightly wrong.

\begin{definition}[LWE ciphertext]
Let $s\in\B^n$ and $\sigma>0$. An LWE encryption of $\mu\in\T$ is the pair
$c=(\mathbf{a},b)\in\T^{n+1}$, where $\mathbf{a}$ is uniform in $\T^n$, $e$ is a
random error drawn from a Gaussian with mean $0$ and standard deviation
$\sigma$, and $b=\langle\mathbf{a},s\rangle+\mu+e$. The \emph{phase} of $c$ is
$\varphi_s(c)=b-\langle\mathbf{a},s\rangle=\mu+e$.
\end{definition}

Because of the error $e$, the equations are only approximately true. Finding
$s$ from them, or even telling them apart from random numbers, is the Learning
With Errors problem, which is believed to be hard even for quantum computers.
So the noise is not a flaw of the scheme: it is what makes it secure. The cost
is that, after removing the mask, we get the phase $\mu+e$ instead of $\mu$.
This is why messages are spread out on the circle and decoded by rounding
(Section~\ref{sec:torus}). The phase is the most important idea in this report:
to know what a ciphertext contains, we compute its phase. For a new ciphertext
with our parameters, $\sigma=2^{-15}\approx3\cdot10^{-5}$, which is more than two
thousand times smaller than the tolerance $\frac1{16}$ of $\Z_8$.

\begin{example}[A small LWE ciphertext]
To keep the numbers small, we use a torus with only 256 points (the integer
$t$ stands for $t/256$), $n=4$, and messages in $\Z_4$, so that the message $1$
is the point $1/4$, i.e., the integer $64$. Take $s=(1,0,1,1)$,
$\mathbf{a}=(100,37,201,54)$ and $e=3$. Then
$\langle\mathbf{a},s\rangle=100+201+54=355$ and $b=355+64+3\equiv166$
(mod $256$). To decrypt, we compute the phase $166-355\equiv67$ and round it to
the closest multiple of $64$, which is $64$: the message is $1$ and the error
was $3$.
\end{example}

\subsection{What LWE Can and Cannot Do}
\label{sec:limits}

\paragraph{Linear operations are free.}
The phase is a linear function of the ciphertext: for integers $k_j$,
$\varphi_s(\sum_jk_jc_j)=\sum_jk_j\varphi_s(c_j)$. So if we add two ciphertexts
component by component, we get a ciphertext of the sum of their messages. For
example, if $c_1$ encrypts $\frac38$ and $c_2$ encrypts $\frac68$, then $c_1+c_2$
decrypts to $\frac98\equiv\frac18$, which is addition modulo $8$. Multiplying a
ciphertext by a public integer works the same way. To add a public constant
$\mu$, we use the \emph{trivial} ciphertext $(\mathbf0,\mu)$, whose phase is
$\mu$ for any key. In short, anyone can compute linear functions of encrypted
messages without knowing the key.

\paragraph{Two problems.}
Linear functions are not enough, and they also have a cost.
\begin{enumerate}
  \item \emph{No nonlinear operation.} To compute any function we need at least
        one nonlinear operation, such as the AND of two bits. LWE ciphertexts
        cannot be multiplied, since their components are points of the torus.
  \item \emph{The noise keeps growing.} When we add two independent
        ciphertexts, the variances of their errors add up, and multiplying by
        $k$ multiplies the standard deviation by $|k|$. After enough
        operations, the error becomes larger than the tolerance and decryption
        fails.
\end{enumerate}

\paragraph{Bootstrapping.}
\cite{gentry2009fhe} showed how to solve the second problem with
\emph{bootstrapping}: running the decryption itself on encrypted data, which
replaces a noisy ciphertext with a fresh ciphertext of the same message. TFHE
uses a version of bootstrapping that solves both problems at once. It takes an
LWE ciphertext with phase $\varphi$, which may be quite noisy, and returns a new
LWE ciphertext of $f(\varphi)$, where $f$ is a function chosen in advance. The
error of the result is small and does \emph{not} depend on the error of the
input. For logic gates, $f$ tells in which half of the circle the phase is.
This function is not linear, which solves the first problem, and its output
has fresh noise, which solves the second.

\subsection{The Idea of TFHE's Bootstrapping: a Table Lookup}
\label{sec:plan}

How can a server compute $f(\varphi)$ if it only has ciphertexts and public keys
and never learns $\varphi$? TFHE uses a table lookup, done by rotating a table,
in three steps (Figure~\ref{fig:plan}):
\begin{enumerate}
  \item round the phase to one of a finite number of positions, $\bar\varphi$;
  \item put in entry $j$ of a table the answer for phases at position $j$,
        that is, $f(j/(2N))$, and \emph{rotate} the table by $\bar\varphi$
        positions, so that the answer for $\bar\varphi$ ends up at position $0$;
  \item read position $0$.
\end{enumerate}
The hard part is step 2: the server must rotate the table by an amount that it
does not know, because $\bar\varphi$ depends on the secret key. To do this, we
need two tools, which are described in the next two subsections:
\begin{itemize}
  \item[(a)] an encryption of a whole table that can be rotated cheaply and
        without increasing its noise: polynomials and TRLWE
        (Section~\ref{sec:trlwe});
  \item[(b)] a way to make each rotation depend on a secret bit, that is, to
        compute ``if the encrypted bit is $1$ then rotate, otherwise do not'':
        the external product and the CMux (Section~\ref{sec:extprod}).
\end{itemize}

% The plan of TFHE's bootstrapping.
\begin{figure}[!htb]
\centering
% The plan of TFHE's bootstrapping: a lookup table stored in a polynomial,
% rotated by the encrypted phase, read at position 0. Shared by report and slides.
\begin{tikzpicture}[font=\footnotesize, >=Stealth]
  % ---------------- encrypted phase ----------------
  \def\rad{1.05}
  \draw[line width=0.9pt] (0,0) circle (\rad);
  \foreach \k in {0,...,15} \draw[cGray] (\k*22.5:\rad-0.06) -- (\k*22.5:\rad+0.06);
  \fill[cGreen] (70:\rad) circle (2.6pt);
  \node[cGreen!75!black, anchor=south] at (70:\rad+0.08) {$\varphi$};
  \node[align=center, font=\scriptsize] at (0,-\rad-0.42) {encrypted phase\\(LWE ciphertext)};
  % ---------------- the table ----------------
  \def\cw{0.66}
  \begin{scope}[shift={(3.2,0.75)}]
    \foreach \j in {0,...,7}
      \node[draw=cNavy!45, fill=cBlue!14, minimum width=\cw cm, minimum height=0.48cm,
            inner sep=0pt] (t\j) at (\j*\cw,0) {$v_{\j}$};
    \node[anchor=west, align=left, font=\scriptsize] at (8*\cw-0.1,0)
      {test polynomial $v$:\\one answer per position};
  \end{scope}
  \begin{scope}[shift={(3.2,-0.85)}, visible on=<3->]
    \foreach \v/\f [count=\j from 0] in {v_3/cBlue!14, v_4/cBlue!14, v_5/cBlue!14,
        v_6/cBlue!14, v_7/cBlue!14, -v_0/cOrange!16, -v_1/cOrange!16, -v_2/cOrange!16}
      \node[draw=cNavy!45, fill=\f, minimum width=\cw cm, minimum height=0.48cm,
            inner sep=0pt] (r\j) at (\j*\cw,0) {$\v$};
    \draw[cGreen!70!black, line width=1.3pt] (-0.5*\cw,-0.24) rectangle (0.5*\cw,0.24);
    \node[anchor=west, align=left, font=\scriptsize] at (8*\cw-0.1,0)
      {$X^{-\bar\varphi}\cdot v$: the answer\\for $\bar\varphi=3$ is now first};
  \end{scope}
  % ---------------- the three steps ----------------
  \draw[->, line width=0.8pt, cNavy!80, visible on=<2->] (1.25,0.45) -- node[above, sloped, font=\scriptsize] {\ding{172} round} (2.75,0.65);
  \node[visible on=<2->, font=\scriptsize, cGray] at (2.05,0.05) {$\bar\varphi=3$};
  \draw[->, line width=0.8pt, cNavy!80, visible on=<3->] (3.2+3.5*\cw,0.48) -- (3.2+3.5*\cw,-0.58)
    node[midway, right, font=\scriptsize, align=left] {\ding{173} rotate by $X^{-\bar\varphi}$\\
      \emph{without} knowing $\bar\varphi$};
  \draw[->, line width=0.8pt, cGreen!70!black, visible on=<4->] (3.2,-1.1) -- (3.2,-1.75)
    node[below, align=center, font=\scriptsize] {\ding{174} read position 0: an\\encryption of $v_{\bar\varphi}$};
\end{tikzpicture}

\caption{The idea of TFHE's bootstrapping, a table lookup by rotation, drawn
with $N=8$. \ding{172} The encrypted phase is rounded to one of $2N=16$
positions. \ding{173} The table of answers, stored as the coefficients of a
polynomial $v$, is rotated by that position, which brings the right answer to
the front (the entries that wrap around come back with a minus sign).
\ding{174} The first entry is read. The hard part is step \ding{173}: the
rotation depends on the secret key, so the server must rotate without knowing
by how much.}
\label{fig:plan}
\end{figure}


\subsection{Polynomials and TRLWE}
\label{sec:trlwe}

\paragraph{Tables as polynomials.}
A table of $N$ points of the torus can be written as a polynomial
$v=v_0+v_1X+\dots+v_{N-1}X^{N-1}$, with one coefficient per entry. Then rotating
the table should be the same as multiplying by a power of $X$. To keep the
degree below $N$, we compute modulo $X^N+1$: we work with $\ZN=\Z[X]/(X^N+1)$
(integer coefficients) and $\TN$ (torus coefficients), where $N$ is a power of
two. The rule $X^N=-1$ turns multiplication by $X$ into a rotation: every
coefficient moves up one position, and the coefficient at the top comes back
at the bottom with its sign changed. No arithmetic is needed, only moving
coefficients and changing signs.

\begin{example}[Rotation with $N=4$]
$(1+2X+3X^3)\cdot X^2=X^2+2X^3+3X^5=-3X+X^2+2X^3$, because
$X^5=X^4\cdot X=-X$. With coefficient vectors, $[1,2,0,3]$ becomes
$[0,-3,1,2]$: every coefficient moves up two positions, and the one that falls
off the end comes back at the start with its sign changed.
\end{example}

Because of the sign change, rotating $N$ times changes the sign of the whole
table ($X^N=-1$), and only after $2N$ rotations do we get back to the start
($X^{2N}=1$). So the exponents of $X$ work modulo $2N$, and a table of $N$
entries has $2N$ rotation positions: in the first $N$ positions we read the
entries $v_j$, and in the next $N$ we read $-v_j$. Section~\ref{sec:testpoly}
shows that this sign change gives us exactly the sign function that the gates
need. Products in $\ZN\times\TN$ are called \emph{negacyclic};
Section~\ref{sec:fft} explains how we compute them quickly.

\paragraph{TRLWE: LWE with polynomials.}
If we replace every number in the definition of LWE by a polynomial, we get an
encryption of a whole table at once.

\begin{definition}[TRLWE ciphertext]
Let $s'\in\BN$ be a polynomial with coefficients in $\{0,1\}$. A TRLWE
encryption of $\mu\in\TN$ is a pair $c=(a,b)\in\TN^2$, where $a$ is uniform,
$e\in\TN$ has independent Gaussian coefficients, and $b=a\cdot s'+\mu+e$. Its
phase is the polynomial $\varphi_{s'}(c)=b-a\cdot s'=\mu+e$.
\end{definition}

Each coefficient of the phase is the matching coefficient of the message plus a
small error. The phase is still linear, so TRLWE ciphertexts can be added just
like LWE ciphertexts. What the plan needs is that they can also be rotated: if
we multiply both components by $X^e$, we get a ciphertext with phase
$\varphi_{s'}(X^ea,X^eb)=X^eb-X^ea\,s'=X^e\varphi_{s'}(c)$. Since $X^e$ only moves
coefficients and changes signs, the error is rotated together with the message
but it does not grow. So rotating an encrypted table by a \emph{known} amount
is free. What is left is to rotate it by a secret amount.

\subsection{Multiplying by an Encrypted Bit}
\label{sec:extprod}

\paragraph{The problem.}
The rotation in the plan depends on the phase $\varphi=b-\sum_ia_is_i$, and
therefore on the secret bits $s_i$. Section~\ref{sec:blindrot} will show that
we can do it one bit at a time, with steps of the form
\[
  \mathrm{ACC}\gets\begin{cases}X^{k_i}\cdot\mathrm{ACC} & \text{if } s_i=1,\\
                                \mathrm{ACC} & \text{if } s_i=0,\end{cases}
\]
which is the same as
\[
  \mathrm{ACC}\gets\mathrm{ACC}+s_i\cdot\big(X^{k_i}-1\big)\cdot\mathrm{ACC},
\]
where $\mathrm{ACC}$ is a TRLWE ciphertext, $k_i$ is public, and the server only
has $s_i$ in encrypted form. So we need to multiply a TRLWE ciphertext by an
encrypted bit, and there are two obstacles. First, the bit cannot be encrypted
with LWE or TRLWE, because then we would have to multiply points of the torus
by points of the torus. Second, even multiplying a ciphertext by a public
integer is risky, because the error is multiplied too: the coefficients of a
ciphertext are arbitrary points of the circle, and written as integers they are
very large numbers. TFHE solves both problems in three steps.

\paragraph{Step 1: small digits.}
The second problem goes away if we never multiply by large numbers, only by
small ones. For this, we write each point of the torus in a base $\Bg$, using a
few digits of small size.

\begin{definition}[Gadget decomposition]
Let $\Bg=2^{\beta_{\mathrm{bit}}}$ and $\ell\ge1$. The decomposition of
$t\in\T$ is a vector $\Dec(t)=(d_1,\dots,d_\ell)\in\Z^\ell$ of signed digits
$d_i\in[-\Bg/2,\Bg/2)$ such that
\begin{equation}
  \Big|\,t-\sum_{i=1}^{\ell}\frac{d_i}{\Bg^{\,i}}\Big|\le\varepsilon
  :=\frac{1}{2\Bg^{\ell}}.
  \label{eq:gadget}
\end{equation}
For a polynomial in $\TN$, we decompose each coefficient, and we get $\ell$
polynomials in $\ZN$ whose coefficients are at most $\beta=\Bg/2$ in absolute
value.
\end{definition}

The vector $(\Bg^{-1},\dots,\Bg^{-\ell})$ that the digits multiply is called the
\emph{gadget}. With our parameters, $\Bg=2^7$ and $\ell=3$, so each coefficient
becomes three digits between $-64$ and $63$, with a rounding error of at most
$2^{-22}$.

\begin{example}[Signed digits in base 10]
\label{ex:decimal}
The idea is easier to see in base $10$ with $\ell=2$. Then the digits must be in
$[-5,5)$ and $\varepsilon=0.005$. For $t=0.7316$, the usual digits are $7$ and
$3$, but $7$ is too large. Since $0.73=1-0.27$, on the torus $0.73$ is the same
point as $-0.27$, so the signed digits are $d_1=-3$ and $d_2=3$. Indeed,
$-\frac3{10}+\frac3{100}=-0.27\equiv0.73$, which is at distance
$0.0016\le\varepsilon$ from $t$. Section~\ref{sec:offset} shows how to compute
such digits quickly.
\end{example}

\paragraph{Step 2: the encryptions that the server needs (TRGSW).}
Let $d=(a,b)$ be a TRLWE ciphertext and $m$ a secret integer (in our case, a
key bit). We want a ciphertext whose phase is
$m\cdot\varphi(d)=m\cdot b-m\,s'\cdot a$. If we decompose the two components
into digit polynomials, $a\approx\sum_iu_i\Bg^{-i}$ and $b\approx\sum_iw_i\Bg^{-i}$,
we get
\begin{equation}
  m\cdot\varphi(d)\;\approx\;\sum_{i=1}^{\ell}w_i\cdot\Big(\frac{m}{\Bg^{\,i}}\Big)
                  +\sum_{i=1}^{\ell}u_i\cdot\Big(\frac{-m\,s'}{\Bg^{\,i}}\Big).
  \label{eq:wish}
\end{equation}
The digits $u_i$ and $w_i$ are small and public, because the server can compute
them from $d$. The terms in parentheses depend on the secret, so they must be
given to the server in encrypted form. A TRGSW ciphertext of $m$ is exactly the
list of these $2\ell$ encryptions.

\begin{definition}[TRGSW ciphertext]
A TRGSW encryption of $m\in\Z$ under $s'$ is a list of $2\ell$ TRLWE encryptions
of zero, $Z_1,\dots,Z_{2\ell}$, to which we add $m$ times the gadget:
$C_i=Z_i+(m/\Bg^{\,i},0)$ and $C_{\ell+i}=Z_{\ell+i}+(0,m/\Bg^{\,i})$ for
$i=1,\dots,\ell$.
\end{definition}

If we compute the phases of the rows, we see that they hold what
Eq.~\eqref{eq:wish} needs: $\varphi(C_i)=e_i-m s'/\Bg^{\,i}$ (the term added to
the first component gets multiplied by $s'$ when we compute the phase) and
$\varphi(C_{\ell+i})=e_{\ell+i}+m/\Bg^{\,i}$ (Figure~\ref{fig:trgsw}). Since each
row is an encryption, the TRGSW ciphertext does not reveal $m$.

% A TRGSW sample as a 2*ell x 2 array and the external product as a weighted
% sum of its rows.
\begin{figure}[!htb]
\centering
% A TRGSW sample as 2*ell TRLWE rows, and the external product as a weighted
% sum of rows by the gadget digits. Shared by the report and the slides.
\begin{tikzpicture}[font=\footnotesize, >=Stealth,
    cell/.style={draw=cNavy!45, minimum height=0.52cm, inner sep=1pt}]
  \def\dy{0.6}
  \node[font=\scriptsize, cGray] at (0,0.62) {digit};
  \node[font=\scriptsize, cGray] at (2.75,0.62) {first component};
  \node[font=\scriptsize, cGray] at (5.75,0.62) {second component};
  \node[font=\scriptsize, cGray, anchor=west] at (7.5,0.62) {phase of the row};
  \foreach \i in {1,2,3} {
    \pgfmathsetmacro{\y}{-(\i-1)*\dy}
    \node[cell, fill=cBlue!14, minimum width=0.75cm] at (0,\y) {$u_{\i}$};
    \node at (0.75,\y) {$\times$};
    \node[cell, fill=white, minimum width=2.7cm] at (2.75,\y)
      {$a_{\i}+{\color{cOrange!90!black}\frac{m}{\Bg^{\i}}}$};
    \node[cell, fill=white, minimum width=2.7cm] at (5.75,\y) {$b_{\i}$};
    \node[anchor=west] at (7.5,\y) {$e_{\i}-{\color{cOrange!90!black}m\,s'/\Bg^{\i}}$};
  }
  \foreach \i in {1,2,3} {
    \pgfmathsetmacro{\y}{-(\i+2)*\dy-0.15}
    \pgfmathtruncatemacro{\r}{\i+3}
    \node[cell, fill=cGreen!14, minimum width=0.75cm] at (0,\y) {$w_{\i}$};
    \node at (0.75,\y) {$\times$};
    \node[cell, fill=white, minimum width=2.7cm] at (2.75,\y) {$a_{\r}$};
    \node[cell, fill=white, minimum width=2.7cm] at (5.75,\y)
      {$b_{\r}+{\color{cOrange!90!black}\frac{m}{\Bg^{\i}}}$};
    \node[anchor=west] at (7.5,\y) {$e_{\r}+{\color{cOrange!90!black}m/\Bg^{\i}}$};
  }
  \draw[decorate, decoration={brace, amplitude=5pt}, cBlue]
    (-0.5,-1.35) -- (-0.5,0.3) node[midway, left=6pt, align=right] {digits\\of $a$};
  \draw[decorate, decoration={brace, amplitude=5pt}, cGreen]
    (-0.5,-3.3) -- (-0.5,-1.65) node[midway, left=6pt, align=right] {digits\\of $b$};
  \draw[line width=0.7pt] (1.3,-3.62) -- (7.2,-3.62);
  \node[anchor=north, align=center] at (4.25,-3.68)
    {$\sum$ of the weighted rows: phase
     $\;m\Big(\underbrace{\textstyle\sum_iw_i\Bg^{-i}}_{\approx b}
       -s'\underbrace{\textstyle\sum_iu_i\Bg^{-i}}_{\approx a}\Big)
       +\underbrace{\textstyle\sum_i(u_ie_i+w_ie_{3+i})}_{\text{small: digits}\,\le\,64}$};
\end{tikzpicture}

\caption{A TRGSW ciphertext of $m$, drawn as $2\ell=6$ TRLWE rows
$(a_r,b_r)$ that encrypt zero, plus the gadget terms (orange). The external
product multiplies each row by one digit polynomial, $u_i$ for the digits of
$a$ and $w_i$ for the digits of $b$, and adds the six rows. Adding the phases on
the right, the gadget terms rebuild $m(b-a\,s')$ from the digits, while the
noise is only multiplied by small digits.}
\label{fig:trgsw}
\end{figure}


\paragraph{Step 3: the external product.}
Now we just combine the rows using the digits as weights, which computes
Eq.~\eqref{eq:wish} on ciphertexts. The \emph{external product} of $C$ with
$d=(a,b)$ is
\begin{equation}
  C\ext d=\sum_{i=1}^{\ell}u_i\cdot C_i+\sum_{i=1}^{\ell}w_i\cdot C_{\ell+i},
  \qquad (u_i)=\Dec(a),\ (w_i)=\Dec(b),
  \label{eq:extprod}
\end{equation}
and, since the phase is linear,
\begin{align*}
\varphi(C\ext d)
  &=m\Big(\sum_i\frac{w_i}{\Bg^{\,i}}-s'\sum_i\frac{u_i}{\Bg^{\,i}}\Big)
    +\sum_i\big(u_ie_i+w_ie_{\ell+i}\big)\\
  &=m\,\varphi(d)+m\,(\epsilon_b-s'\epsilon_a)
    +\sum_i\big(u_ie_i+w_ie_{\ell+i}\big),
\end{align*}
where $\epsilon_a$ and $\epsilon_b$ are the decomposition errors, whose
coefficients are at most $\varepsilon$. Both problems are solved: we never
multiply two points of the torus (the digits are integers), and the errors
$e_r$ of the rows are only multiplied by digits of size at most $\beta=64$.

\begin{proposition}[Noise of the external product {\cite{chillotti2020tfhe}}]
\label{prop:extprod}
Let $C$ encrypt $m\in\{0,1\}$ with noise variance $\vartheta_C$, and let the error
of $d$ have variance $\vartheta_d$. Then the error of $C\ext d$ has variance at
most $2\ell N\beta^2\vartheta_C+m\big((1+N)\varepsilon^2+\vartheta_d\big)$.
\end{proposition}

The three terms come from the last equation. The first one comes from the
$2\ell$ products of a digit polynomial by an error polynomial: each coefficient
of such a product is a sum of $N$ terms, each a digit (at most $\beta$) times an
error. The second comes from the decomposition errors, and the third is the
error that $d$ already had. The first term explains how the parameters are
chosen. A larger base $\Bg$ needs fewer digits $\ell$, so there are fewer
products and the external product is faster, but $\beta$ grows and so does the
noise. It also explains why TRGSW ciphertexts need a very small error: with our
parameters $2\ell N\beta^2\approx2.5\cdot10^7$, so they are created with standard
deviation $2^{-25}$ instead of the $2^{-15}$ used for LWE
(Table~\ref{tab:params}). Finally, the error of $d$ is carried over but not
multiplied, so the noise grows by \emph{adding} a small amount each time. This
is what allows hundreds of external products in a row.

\paragraph{The encrypted if-then-else (CMux).}
With multiplication by an encrypted bit, we can now choose between two
ciphertexts without knowing the choice. Given TRLWE ciphertexts $d_0$ and $d_1$
and a TRGSW ciphertext $C$ of a bit $m$,
\begin{equation}
  \mathrm{CMux}(C,d_1,d_0)=d_0+C\ext(d_1-d_0)
  \label{eq:cmux}
\end{equation}
encrypts $\mu_0+m(\mu_1-\mu_0)$, that is, the message of $d_1$ if $m=1$ and the
message of $d_0$ if $m=0$, and the server does not learn which one
(Figure~\ref{fig:cmux}a). It costs one external product in both cases. If we
take $d_1=X^{k_i}\cdot\mathrm{ACC}$, $d_0=\mathrm{ACC}$ and $C$ a TRGSW
ciphertext of $s_i$, this is exactly the step that we needed at the start of
this subsection.

\begin{example}[CMux]
\label{ex:cmux}
Let $d_0$ and $d_1$ encrypt the constant polynomials $\mu_0=\frac18$ and
$\mu_1=\frac58$. The difference $d_1-d_0$ encrypts $\frac12$, so $C\ext(d_1-d_0)$
encrypts $0$ if $m=0$ and $\frac12$ if $m=1$. Adding $d_0$ gives a phase of
$\frac18$ or $\frac18+\frac12=\frac58$, which is the message of $d_0$ or of $d_1$
(Figure~\ref{fig:cmux}b).
\end{example}

% CMux as an encrypted multiplexer, and the example used by test06 and test07.
\begin{figure}[!htb]
\centering
% CMux: (a) the encrypted multiplexer; (b) the example d0 = 1/8, d1 = 5/8.
% Shared by the report and the slides.
\begin{tikzpicture}[font=\footnotesize, >=Stealth]
  % ---------------- (a) ----------------
  \draw[line width=0.8pt, fill=cBlue!10] (0,-1.15) -- (0,1.15) -- (1.4,0.62) -- (1.4,-0.62) -- cycle;
  \node[font=\small\bfseries] at (0.66,0) {CMux};
  \node[font=\scriptsize, cGray] at (0.2,0.7) {1};
  \node[font=\scriptsize, cGray] at (0.2,-0.7) {0};
  \draw[->, line width=0.8pt] (-1.2,0.7) node[left] {$d_1$} -- (0,0.7);
  \draw[->, line width=0.8pt] (-1.2,-0.7) node[left] {$d_0$} -- (0,-0.7);
  \draw[->, line width=0.8pt, cOrange] (0.7,1.85) node[above] {$C=\mathrm{TRGSW}(m)$} -- (0.7,0.9);
  \draw[->, line width=0.8pt] (1.4,0) -- (2.3,0);
  \node[anchor=west, align=left] at (2.35,0) {$d_0+C\ext(d_1-d_0)$};
  \node[anchor=north west, align=left, cGray!90!black, font=\scriptsize] at (2.35,-0.3)
    {message $\mu_0+m(\mu_1-\mu_0)$:\\$\mu_1$ if $m=1$, \ $\mu_0$ if $m=0$};
  \node at (-1.9,1.9) {(a)};
  % ---------------- (b) ----------------
  \begin{scope}[xshift=9.0cm]
    \def\rad{1.45}
    \draw[line width=0.9pt] (0,0) circle (\rad);
    \foreach \k in {0,...,7} \draw[cGray] (\k*45:\rad-0.07) -- (\k*45:\rad+0.07);
    \draw[->, cOrange, line width=1.1pt, visible on=<3->] (52:\rad-0.25) arc (52:218:\rad-0.25);
    \fill[cBlue] (45:\rad) circle (2.6pt);
    \fill[cOrange] (225:\rad) circle (2.6pt);
    \node[cBlue!85!black, anchor=south west] at (45:\rad) {$d_0$: $\frac18$};
    \node[cOrange!85!black, anchor=north east] at (225:\rad) {$d_1$: $\frac58$};
    \node[cOrange!85!black, align=center, visible on=<3->] at (135:0.55)
      {$m=1$\\$+\frac12$};
    \node[cBlue!85!black, anchor=west, align=left, visible on=<2->] at (5:\rad+0.2)
      {$m=0$: $+0$,\\stays at $\frac18$};
    \node[cGray!90!black, align=center, font=\scriptsize] at (0,-\rad-0.55)
      {$d_1-d_0$ encrypts $\frac12$};
    \node at (-\rad-1.0,1.9) {(b)};
  \end{scope}
\end{tikzpicture}

\caption{The CMux. (a) It chooses one of two TRLWE ciphertexts according to
a bit $m$ that is itself encrypted, using one external product. (b) The case of
Example~\ref{ex:cmux}: $d_0$ encrypts $\frac18$ and $d_1$ encrypts $\frac58$,
so $d_1-d_0$ encrypts $\frac12$. The external product adds $m\cdot\frac12$ to the
phase of $d_0$, which stays at $\frac18$ when $m=0$ and moves half a turn, to
$\frac58$, when $m=1$. The server does the same work in both cases.}
\label{fig:cmux}
\end{figure}


%% ===========================================================================
\section{How a Gate Is Bootstrapped}
\label{sec:bootstrapping}

\subsection{Overview}

We now have everything we need to carry out the plan of
Section~\ref{sec:plan} on real ciphertexts. Figure~\ref{fig:pipeline} shows the
path of one gate. Each stage has a clear purpose:
\begin{itemize}
  \item the \emph{linear combination} of the two inputs gives a ciphertext whose
        phase is in the upper half of the circle exactly when the output of the
        gate should be \textsf{true} (Section~\ref{sec:gates});
  \item \emph{modulus switching} turns its phase, a point of the torus, into one
        of $2N$ rotation positions (Section~\ref{sec:modswitch});
  \item \emph{blind rotation} rotates the \emph{test polynomial}, which is the
        table of answers (Section~\ref{sec:testpoly}), by that position, using
        one CMux per bit of the secret key (Section~\ref{sec:blindrot});
  \item \emph{sample extraction} reads position $0$ of the rotated table as an
        LWE ciphertext (Section~\ref{sec:extract});
  \item \emph{key switching} brings this ciphertext back to the key and size of
        the inputs, so that it can be used by the next gate
        (Section~\ref{sec:ks}).
\end{itemize}

\paragraph{Who has which key.}
The client has two secret keys: $s\in\B^n$, used to encrypt and decrypt the
user's bits as LWE ciphertexts, and $s'\in\BN$, used for the TRLWE ciphertexts
inside the bootstrapping. The server must run CMux operations controlled by
the bits of $s$, so it receives each of these bits encrypted as a TRGSW
ciphertext under $s'$. This is the \emph{bootstrapping key},
$\BK_i=\mathrm{TRGSW}_{s'}(s_i)$ for $i=1,\dots,n$. The server also receives the
\emph{key-switching key}, described in Section~\ref{sec:ks}. Together, they form
the public \emph{evaluation key}. Publishing it is considered safe under the
usual assumption, called circular security, that each key can be encrypted
under the other one.

% The four stages of a bootstrapped gate and the two parts of the evaluation key.
\begin{figure}[!htb]
\centering
\begin{tikzpicture}[
    font=\footnotesize, >=stealth,
    stage/.style={draw, rounded corners=2pt, line width=0.7pt, fill=#1,
                  minimum height=1.05cm, text width=1.75cm, align=center,
                  inner sep=2pt},
    stage/.default=cBlue!10,
    kbox/.style={draw, densely dashed, rounded corners=2pt, fill=cGray!8,
                text width=3.3cm, align=center, font=\footnotesize,
                inner sep=3pt},
    lbl/.style={font=\scriptsize, align=center, cGray!90!black},
    flow/.style={->, line width=0.8pt}]
  \node[stage=cGray!6, text width=0.9cm] (in) {$c_1, c_2$};
  \node[stage, right=1.0cm of in]  (lin) {linear\\combination};
  \node[stage, right=1.0cm of lin] (br)  {blind\\rotation};
  \node[stage, right=1.0cm of br]  (se)  {sample\\extraction};
  \node[stage, right=1.0cm of se]  (ks)  {key\\switching};
  \node[stage=cGray!6, text width=0.9cm, right=1.0cm of ks] (out) {$c$};
  \draw[flow] (in)  -- node[lbl, above] {LWE} node[lbl, below] {$s,\,n$} (lin);
  \draw[flow] (lin) -- node[lbl, above] {LWE} node[lbl, below] {$s,\,n$} (br);
  \draw[flow] (br)  -- node[lbl, above] {TRLWE} node[lbl, below] {$s',\,N$} (se);
  \draw[flow] (se)  -- node[lbl, above] {LWE} node[lbl, below] {$s',\,N$} (ks);
  \draw[flow] (ks)  -- node[lbl, above] {LWE} node[lbl, below] {$s,\,n$} (out);
  \node[kbox, below=0.75cm of br] (bk) {bootstrapping key\\$n$ TRGSW ciphertexts of $s_i$ under $s'$};
  \node[kbox, below=0.75cm of ks] (ksk) {key-switching key\\$N t\,B_{\mathrm{ks}}$ LWE ciphertexts under $s$};
  \draw[->, densely dashed, cGray] (bk) -- (br);
  \draw[->, densely dashed, cGray] (ksk) -- (ks);
  \node[lbl, above=1pt of lin] {cheap, adds noise};
  \node[lbl, above=1pt of br] {$>95\%$ of the time};
  \node[lbl, above=1pt of se] {free, no noise};
  \node[lbl, above=1pt of ks] {changes key};
\end{tikzpicture}
\caption{The path of a bootstrapped gate. The two input ciphertexts are
combined linearly; blind rotation computes the gate on the encrypted phase and
gives a TRLWE ciphertext with fresh noise; sample extraction reads its constant
coefficient as an LWE ciphertext under the key $s'$, of size $N$; and key
switching brings it back to the key $s$, of size $n$, so that it can be used by
the next gate. The dashed boxes are the two parts of the public evaluation key
that the server has.}
\label{fig:pipeline}
\end{figure}


\subsection{Step 1: Modulus Switching}
\label{sec:modswitch}

We can only rotate by whole numbers of positions, modulo $2N$
(Section~\ref{sec:trlwe}), but the phase $\varphi=b-\sum_ia_is_i$ of the input
$c=(\mathbf a,b)$ is a point of the torus. So the first step cuts the circle
into $2N$ slices and asks in which slice the phase is. The server cannot
compute $\varphi$, so instead it rounds every public component of $c$:
\[
  \bar b=\round{2Nb},\qquad \bar a_i=\round{2Na_i},\qquad
  \bar\varphi=\bar b-\sum_i\bar a_is_i\bmod 2N\approx2N\varphi.
\]
The integer $\bar\varphi$ is the number of the slice
(Figure~\ref{fig:testpoly}a shows the slices for $N=4$). It is still secret,
because it depends on $s$, but now it is a valid rotation amount. Rounding adds
a small error: each rounded value is off by at most $1/(4N)$, with variance
$1/(48N^2)$, and only $\bar b$ and the $\bar a_i$ with $s_i=1$ (about half of
them) matter. This \emph{modulus-switching error} has variance
$\vartheta_{\mathrm{ms}}\approx(1+n/2)/(48N^2)\approx6.3\cdot10^{-6}$, a standard
deviation of $0.0025$. It adds to the error of the input and can at most make
the lookup use a neighboring slice; it never reaches the output.

\subsection{Step 2: The Test Polynomial}
\label{sec:testpoly}

After modulus switching, the input has been reduced to a slice number
$\bar\varphi$ between $0$ and $2N-1$. The bootstrapping must now return a value
that depends on this slice, and the \emph{test polynomial} is the table that
says which value.

\paragraph{What the coefficients $v_j$ are.}
The test polynomial $v=v_0+v_1X+\dots+v_{N-1}X^{N-1}$ has torus coefficients
chosen by the server. Each coefficient is one entry of the table:
\begin{quote}
  $v_j$ is the value that the bootstrapping returns when the rounded phase of
  its input is $\bar\varphi=j$, that is, when the phase is close to $j/(2N)$.
\end{quote}
The $N$ coefficients cover the $N$ slices of the upper half of the circle. The
slices of the lower half do not have their own coefficients: slice $N+j$, which
is opposite slice $j$, returns $-v_j$. Figure~\ref{fig:testwheel} shows this for
$N=4$, with each of the $8$ slices labeled with its output.

So, choosing $v$ means choosing what the bootstrapping computes: to evaluate a
function $f$ of the phase, we set $v_j=f\big(j/(2N)\big)$ for $j=0,\dots,N-1$.
Figure~\ref{fig:testwheel} shows two examples.
\begin{itemize}
  \item \emph{Gates} (b): $v_j=\mu=\frac18$ for every $j$. Every slice of the
        upper half returns $+\frac18$ and every slice of the lower half
        returns $-\frac18$. The bootstrapping computes the sign of the phase,
        already in the encoding of bits (\textsf{true} is $+\frac18$,
        \textsf{false} is $-\frac18$). The sign change caused by $X^N=-1$ gives
        us the second half of the sign function for free.
  \item \emph{A table that is not constant} (c): with $N=4$ and $v_j=j/8$, slice
        $j$ of the upper half returns $j/8$. An input whose phase is close to
        $\frac28$ comes out as an encryption of $\frac28$, with fresh noise.
\end{itemize}
Since the lower half always returns the negated entries, the function must
satisfy $f(\varphi+\frac12)=-f(\varphi)$. The sign function does. Other
functions can still be computed if we make sure that the inputs are always in
the upper half of the circle, as in the example at the end of this subsection.

% What the coefficients of the test polynomial mean.
\begin{figure}[!htb]
\centering
% What the coefficients v_j of the test polynomial mean, for N = 4: each of the
% 2N slices of the torus is labeled with the output that the bootstrapping
% returns for phases in that slice. (a) any v; (b) the gate test polynomial;
% (c) v_j = j/8. Shared by the report and the slides.
\begin{tikzpicture}[font=\footnotesize]
  \def\rad{1.42}
  % \wheel{x shift}{labels of the 8 slices}{title}{polynomial}{overlay}
  \def\wheel#1#2#3#4#5{%
    \begin{scope}[xshift=#1cm, visible on=#5]
      \foreach \lab [count=\k from 0] in {#2} {
        \pgfmathsetmacro{\aa}{\k*45-22.5}
        \pgfmathsetmacro{\bb}{\k*45+22.5}
        \ifnum\k<4
          \fill[cBlue!14] (0,0) -- (\aa:\rad) arc (\aa:\bb:\rad) -- cycle;
        \else
          \fill[cOrange!16] (0,0) -- (\aa:\rad) arc (\aa:\bb:\rad) -- cycle;
        \fi
        \draw[white, line width=1pt] (0,0) -- (\aa:\rad);
        \node at (\k*45:0.66*\rad) {\lab};
        \node[font=\scriptsize, cGray] at (\k*45:\rad+0.24) {\k};
      }
      \draw[cNavy, line width=0.9pt] (0,0) circle (\rad);
      \fill[cNavy] (0,0) circle (0.8pt);
      \node[anchor=north, align=center] at (0,-\rad-0.38) {#3\\[1pt]#4};
    \end{scope}}
  \wheel{0}{$v_0$,$v_1$,$v_2$,$v_3$,$-v_0$,$-v_1$,$-v_2$,$-v_3$}
        {\textbf{(a)} any test polynomial}{$v=v_0+v_1X+v_2X^2+v_3X^3$}{<1->}
  \wheel{4.75}{$+\frac18$,$+\frac18$,$+\frac18$,$+\frac18$,$-\frac18$,$-\frac18$,$-\frac18$,$-\frac18$}
        {\textbf{(b)} gates: all $v_j=\frac18$}{$v=\frac18(1+X+X^2+X^3)$}{<2->}
  \wheel{9.5}{$0$,$\frac18$,$\frac14$,$\frac38$,$0$,$-\frac18$,$-\frac14$,$-\frac38$}
        {\textbf{(c)} $v_j=j/8$}{$v=\frac18X+\frac14X^2+\frac38X^3$}{<3->}
\end{tikzpicture}

\caption{What the coefficients $v_j$ of the test polynomial mean, for $N=4$.
The circle is cut into $2N=8$ slices, numbered by the rounded phase
$\bar\varphi$ (outside), and each slice shows the value that the bootstrapping
returns when the phase of its input is in that slice. Coefficient $v_j$ is the
output for slice $j$ of the upper half (blue); the opposite slice $N+j$, in the
lower half (orange), always returns $-v_j$, because $X^N=-1$. (a) The general
rule. (b) The gate test polynomial: every $v_j$ is $\frac18$, so the output is
$+\frac18$ on the upper half and $-\frac18$ on the lower half. (c) A table that
is not constant, $v_j=j/8$: on the upper half the bootstrapping returns the
rounded phase itself.}
\label{fig:testwheel}
\end{figure}


\paragraph{Why rotating reads the right entry.}
Now we can see why multiplying by $X^{-\bar\varphi}$ looks up entry
$\bar\varphi$. Multiplying by $X^{-1}$ moves every coefficient one position
toward the constant term: the coefficient of $X$ becomes the constant
coefficient, and the old constant coefficient, which falls off the end, comes
back as the coefficient of $X^{N-1}$ with its sign changed. After $\bar\varphi$
such moves, the constant coefficient is $v_{\bar\varphi}$ if $\bar\varphi<N$, and
$-v_{\bar\varphi-N}$ otherwise, which is exactly the label of slice $\bar\varphi$
in Figure~\ref{fig:testwheel}. Figure~\ref{fig:testpoly}b lists all $2N$
rotations for $N=4$. Reading its first column from top to bottom gives
$v_0,\dots,v_3,-v_0,\dots,-v_3$, the labels of the slices in order around the
circle.

% The test polynomial as a lookup table, for N = 4.
\begin{figure}[!htb]
\centering
% The test polynomial as a lookup table, for N = 4: (a) the 2N positions of the
% rounded phase on the torus; (b) X^(-phibar) * v for every phibar, whose
% constant coefficient (framed) is the output. Shared by the report and slides.
\begin{tikzpicture}[font=\footnotesize, >=Stealth]
  % ---------------- (a) ----------------
  \begin{scope}
    \def\rad{1.5}
    \fill[cBlue!12] (0,0) -- (\rad,0) arc (0:180:\rad) -- cycle;
    \fill[cOrange!14] (0,0) -- (-\rad,0) arc (180:360:\rad) -- cycle;
    \draw[line width=0.9pt] (0,0) circle (\rad);
    \foreach \p in {0,...,7} {
      \fill (\p*45:\rad) circle (1.8pt);
      \node at (\p*45:\rad+0.38) {$\p$};
    }
    \fill[cGreen] (51:\rad) circle (2.2pt);
    \draw[->, cGreen, line width=0.9pt] (51:\rad-0.05) arc (51:46:\rad-0.05);
    \node[cGreen!80!black, align=center, font=\scriptsize] at (0,0.62)
      {phase $\varphi$ $\mapsto$ $\bar\varphi=1$};
    \node[cBlue!85!black, font=\scriptsize] at (0,0.2) {$\bar\varphi<N$};
    \node[cOrange!85!black, font=\scriptsize] at (0,-0.3) {$\bar\varphi\ge N$};
    \node[align=center, font=\scriptsize, cGray] at (0,-\rad-0.8) {$\bar\varphi=\round{2N\varphi}$, \ $2N=8$};
    \node at (-\rad-0.6,\rad+0.35) {(a)};
  \end{scope}
  % ---------------- (b) ----------------
  \begin{scope}[xshift=4.3cm, yshift=1.65cm]
    \def\cw{0.82}
    \def\ch{0.42}
    \node[font=\scriptsize, cGray] at (-0.55,0.42) {$\bar\varphi$};
    \foreach \j in {0,...,3}
      \node[font=\scriptsize, cGray] at (\j*\cw,0.42) {$X^{\j}$};
    \foreach \r in {0,...,7} {
      \node[font=\scriptsize] at (-0.55,-\r*\ch) {\r};
      \foreach \j in {0,...,3} {
        \pgfmathtruncatemacro{\k}{\j+\r}
        \pgfmathtruncatemacro{\q}{mod(floor(\k/4),2)}
        \pgfmathtruncatemacro{\c}{mod(\k,4)}
        \ifnum\q=0
          \node[draw=cNavy!35, fill=cBlue!14, minimum width=\cw cm, minimum height=\ch cm,
                inner sep=0pt] at (\j*\cw,-\r*\ch) {$v_{\c}$};
        \else
          \node[draw=cNavy!35, fill=cOrange!16, minimum width=\cw cm, minimum height=\ch cm,
                inner sep=0pt] at (\j*\cw,-\r*\ch) {$-v_{\c}$};
        \fi
      }
    }
    \draw[cGreen!70!black, line width=1.3pt] (-0.5*\cw,0.5*\ch) rectangle (0.5*\cw,-7.5*\ch);
    \node[cGreen!70!black, font=\scriptsize, anchor=north west] at (-0.5*\cw,-7.6*\ch)
      {$\uparrow$ constant coefficient};
    \node[font=\scriptsize, align=center, cGray] at (1.5*\cw,-9.9*\ch)
      {$X^{-\bar\varphi}\cdot v$, \ $v=v_0+v_1X+v_2X^2+v_3X^3$};
    % output with the gate test polynomial
    \begin{scope}[visible on=<2->]
      \node[font=\scriptsize, cGray, align=center] at (5.0*\cw,0.42) {gates: $v_j=\mu$};
      \foreach \r in {0,...,3}
        \node[cBlue!85!black] at (5.0*\cw,-\r*\ch) {$+\mu$};
      \foreach \r in {4,...,7}
        \node[cOrange!85!black] at (5.0*\cw,-\r*\ch) {$-\mu$};
      \draw[->, cGray] (3.65*\cw,-1.5*\ch) -- (4.45*\cw,-1.5*\ch);
      \draw[->, cGray] (3.65*\cw,-5.5*\ch) -- (4.45*\cw,-5.5*\ch);
    \end{scope}
    \node at (-1.05,0.6) {(b)};
  \end{scope}
\end{tikzpicture}

\caption{The test polynomial as a lookup table, drawn for $N=4$. (a) The
rounded phase $\bar\varphi=\round{2N\varphi}$ is one of $2N=8$ positions on the
torus; positions $0$ to $N-1$ are the upper half and $N$ to $2N-1$ the lower
half. (b) The coefficients of $X^{-\bar\varphi}\cdot v$ for every $\bar\varphi$:
each row is the row above shifted to the left, and the coefficient that leaves
on the left comes back on the right with its sign changed, because $X^N=-1$.
The constant coefficient (framed) is $v_{\bar\varphi}$ in the upper half and
$-v_{\bar\varphi-N}$ in the lower half. With every $v_j=\mu$, the output is $+\mu$
or $-\mu$ depending on the half of the circle where the phase is: the sign
function used by every gate.}
\label{fig:testpoly}
\end{figure}


\paragraph{Is the test polynomial encrypted?}
No. The function being computed is not a secret (the server knows that it is
computing a NAND, for example), so the table is public. What must stay hidden
is the input, and therefore \emph{which} entry of the table is selected. The
test polynomial enters the blind rotation as the \emph{trivial} TRLWE
ciphertext $(0,X^{-\bar b}\cdot v)$. Its mask is zero, so its phase is
$X^{-\bar b}v$ under any key, without any noise. This ciphertext hides nothing,
but it does not need to, since $\bar b$ comes from the public part of the
input. The accumulator becomes a real ciphertext at the first CMux, which adds
$\BK_i\ext\big((X^{\bar a_i}-1)\cdot\mathrm{ACC}\big)$. This term is a new
encryption, under the real key $s'$, of either the rotation difference (if
$s_i=1$) or zero (if $s_i=0$). From then on the accumulator has a random mask and
some noise, and the two possible results look the same to anyone without the
key (Figure~\ref{fig:trivial}). In other words, the table is public, but the
choice of the entry is secret.

% The test polynomial is public: it enters blind rotation as a trivial sample.
\begin{figure}[!htb]
\centering
% The test polynomial enters blind rotation as a trivial (public, noiseless)
% TRLWE sample; the first CMux turns the accumulator into a real ciphertext.
% Example of Section 3 (N = 4, b-bar = 6). Shared by the report and the slides.
\begin{tikzpicture}[font=\footnotesize, >=Stealth]
  \def\cw{0.78}
  \def\ch{0.48}
  % ---------------- trivial sample ----------------
  \node[font=\small\bfseries, anchor=south] at (1.5*\cw,0.62) {start: trivial ciphertext};
  \node[anchor=east] at (-0.5*\cw-0.1,0.2) {mask $a$};
  \node[anchor=east] at (-0.5*\cw-0.1,-0.35) {body $b$};
  \foreach \j in {0,...,3}
    \node[draw=cNavy!40, fill=white, minimum width=\cw cm, minimum height=\ch cm,
          inner sep=0pt] at (\j*\cw,0.2) {$0$};
  \foreach \v/\f [count=\j from 0] in {-\mu/cOrange!16, -\mu/cOrange!16, +\mu/cBlue!14, +\mu/cBlue!14}
    \node[draw=cNavy!40, fill=\f, minimum width=\cw cm, minimum height=\ch cm,
          inner sep=0pt] at (\j*\cw,-0.35) {$\v$};
  \node[anchor=north, align=center, font=\scriptsize] at (1.5*\cw,-0.7)
    {$=X^{-6}v$, \ phase $=b$ under \emph{any} key\\public, no noise};
  % ---------------- the CMux ----------------
  \draw[->, line width=0.9pt, cNavy!80, visible on=<2->] (3.9*\cw,-0.07) -- (7.1*\cw,-0.07)
    node[midway, above, align=center, font=\scriptsize] {CMux with $\BK_1$}
    node[midway, below, align=center, font=\scriptsize, cGray!90!black]
      {adds an encryption\\under the real key};
  % ---------------- real ciphertext ----------------
  \begin{scope}[xshift=8.1*\cw cm, visible on=<2->]
    \node[font=\small\bfseries, anchor=south] at (1.5*\cw,0.62) {after step 1: real ciphertext};
    \foreach \j in {0,...,3}
      \node[draw=cNavy!40, fill=cGray!22, minimum width=\cw cm, minimum height=\ch cm,
            inner sep=0pt, text=cGray!60!black] at (\j*\cw,0.2) {$?$};
    \foreach \j in {0,...,3}
      \node[draw=cNavy!40, fill=cGray!22, minimum width=\cw cm, minimum height=\ch cm,
            inner sep=0pt, text=cGray!60!black] at (\j*\cw,-0.35) {$?$};
    \node[anchor=north, align=center, font=\scriptsize] at (1.5*\cw,-0.7)
      {random-looking without the key;\\phase $=b-a\cdot s'=X^{-3}v+e$};
  \end{scope}
\end{tikzpicture}

\caption{The test polynomial is not encrypted. Blind rotation starts from the
trivial TRLWE ciphertext $(0,X^{-\bar b}v)$. Its mask is zero, so its phase is
the public polynomial $X^{-\bar b}v$ under any key (here $N=4$, $\bar b=6$ and
the gate test polynomial, as in Example~\ref{ex:blindrot}). The first CMux adds
an external product with a ciphertext of the bootstrapping key, which is a
real encryption under $s'$. After it, the accumulator has a random mask and
some noise, and without the key one cannot see whether it was rotated (here by
$X^{3}$).}
\label{fig:trivial}
\end{figure}


\paragraph{Programmable bootstrapping.}
Using the bootstrapping to compute functions other than the sign was later
called \emph{programmable bootstrapping} \cite{chillotti2021pbs}. For example,
with $N=1024$ we can split the coefficients into four bands of $N/4=256$
consecutive coefficients and store $m/8$ in every coefficient of band $m$, that
is, $v_j=\lfloor4j/N\rfloor/8$. An input encrypting $(2m+1)/16$ has a rounded
phase close to $2N\cdot(2m+1)/16=128\,(2m+1)$, which is the middle of band $m$,
so it comes out as an encryption of $m/8$. In this way one bootstrapping
computes a lookup table with four entries instead of a sign. All these inputs
are in the upper half of the circle, so the negated entries of the lower half
are never used.

\subsection{Step 3: Blind Rotation}
\label{sec:blindrot}

We still need to compute an encryption of $X^{-\bar\varphi}\cdot v$ without
knowing $\bar\varphi$. The main observation is that the rotation can be split
into one factor per key bit:
\[
  X^{-\bar\varphi}=X^{-\bar b}\cdot\prod_{i=1}^{n}X^{\bar a_is_i},
  \qquad
  X^{\bar a_is_i}=\begin{cases}X^{\bar a_i}&\text{if }s_i=1,\\1&\text{if }s_i=0.\end{cases}
\]
The first factor only uses public data. Each of the other factors is an
if-then-else that depends on a secret bit, which is exactly what a CMux with
$\BK_i=\mathrm{TRGSW}(s_i)$ computes. So the server starts from the trivial
TRLWE ciphertext $\mathrm{ACC}=(0,X^{-\bar b}v)$ and, for $i=1,\dots,n$, replaces
it by
\[
  \mathrm{CMux}\big(\BK_i,\;X^{\bar a_i}\cdot\mathrm{ACC},\;\mathrm{ACC}\big)
  =\mathrm{ACC}+\BK_i\ext\big((X^{\bar a_i}-1)\cdot\mathrm{ACC}\big).
\]
After the $n$ steps, the accumulator encrypts $X^{-\bar\varphi}\cdot v$. The
server knows every $\bar a_i$ and $\bar b$, but it does not know which of the
$n$ steps actually rotated, so it never learns the total rotation. This is why
the operation is called \emph{blind rotation}. Each step costs one external
product and adds its error (Proposition~\ref{prop:extprod}), which is why the
bootstrapping key must be so precise.

\begin{example}[Bootstrapping an encryption of \textsf{true}, step by step]
\label{ex:blindrot}
Take $n=3$, $N=4$ (so there are $2N=8$ positions), the secret key $s=(1,0,1)$,
and the gate test polynomial $v=\mu\,(1+X+X^2+X^3)$ with $\mu=\frac18$.

\emph{The input.} Let $c$ be an LWE encryption of \textsf{true}, that is, of
$+\frac18$, with error $e=0.003$. With the mask $\mathbf a=(0.37,\,0.61,\,0.24)$,
the body is $b=0.37\cdot1+0.61\cdot0+0.24\cdot1+0.125+0.003=0.738$, and the
phase is $b-(0.37+0.24)=0.128$.

\emph{Step 1, modulus switching.} We multiply by $2N=8$ and round:
$\bar b=\round{5.904}=6$ and
$\bar{\mathbf a}=(\round{2.96},\round{4.88},\round{1.92})=(3,5,2)$. So
$\bar\varphi=6-(3\cdot1+5\cdot0+2\cdot1)=1$, a position in the upper half, as
expected for a phase close to $\frac18$. The server knows $\bar b$ and
$\bar{\mathbf a}$, but not $\bar\varphi$.

\emph{Step 3, blind rotation.} The accumulator starts as the trivial ciphertext
$(0,X^{-6}v)$. Since $X^8=1$, we have $X^{-6}=X^{2}$, and
$X^2v=\mu(X^2+X^3+X^4+X^5)=\mu(-1-X+X^2+X^3)$, because $X^4=-1$ and $X^5=-X$.
The server then runs one CMux per key bit. It always does the same work, and
the key decides which steps rotate:
\begin{center}
\small
\begin{tabular}{@{}lccccc@{}}
\toprule
step & $\bar a_i$ & $s_i$ & effect & message of ACC & coefficients\\
\midrule
start & -- & -- & trivial ciphertext & $X^{-6}v$ & $[-\mu,-\mu,+\mu,+\mu]$\\
$i=1$ & $3$ & $1$ & $\times X^{3}$ & $X^{-3}v$ & $[+\mu,-\mu,-\mu,-\mu]$\\
$i=2$ & $5$ & $0$ & unchanged & $X^{-3}v$ & $[+\mu,-\mu,-\mu,-\mu]$\\
$i=3$ & $2$ & $1$ & $\times X^{2}$ & $X^{-1}v$ & $[+\mu,+\mu,+\mu,-\mu]$\\
\bottomrule
\end{tabular}
\end{center}
From step $i=1$ on, the accumulator is a real ciphertext
(Figure~\ref{fig:trivial}). The table shows the \emph{messages} it encrypts,
which the server never sees.

\emph{Step 4, sample extraction.} The constant coefficient of the final message
is $+\mu=+\frac18$, so the output is an encryption of \textsf{true}, the sign of
the input phase, with fresh noise. Figure~\ref{fig:blindrot} shows the same
three steps for any test polynomial; there the result is $v_1=v_{\bar\varphi}$,
the entry for slice $1$ (phases close to $\frac18$). With the gate test
polynomial every $v_j$ is $\mu$, so any $\bar\varphi\in\{0,1,2,3\}$ would give
$+\mu$ and any $\bar\varphi\in\{4,\dots,7\}$ would give $-\mu$.
\end{example}

% Blind rotation step by step, for n = 3 and N = 4.
\begin{figure}[!htb]
\centering
% Blind rotation step by step, for n = 3 and N = 4. Shared by the report and
% the slides (each step appears on its own overlay in the slides).
\begin{tikzpicture}[font=\footnotesize, >=Stealth]
  \def\cw{0.82}
  \def\ch{0.46}
  \def\dy{1.05}
  % one accumulator: \acc{row}{sign/index list}{overlay}; b = +v_index, o = -v_index
  \def\acc#1#2#3{%
    \begin{scope}[shift={(0,-#1*\dy)}, visible on=#3]
      \foreach \s/\c [count=\j from 0] in {#2} {
        \if\s b
          \node[draw=cNavy!35, fill=cBlue!14, minimum width=\cw cm, minimum height=\ch cm,
                inner sep=0pt] at (\j*\cw,0) {$v_{\c}$};
        \else
          \node[draw=cNavy!35, fill=cOrange!16, minimum width=\cw cm, minimum height=\ch cm,
                inner sep=0pt] at (\j*\cw,0) {$-v_{\c}$};
        \fi
      }
      \draw[cGreen!70!black, line width=1.2pt] (-0.5*\cw,-0.5*\ch) rectangle (0.5*\cw,0.5*\ch);
    \end{scope}}
  \acc{0}{o/2, o/3, b/0, b/1}{<1->}
  \acc{1}{b/3, o/0, o/1, o/2}{<2->}
  \acc{2}{b/3, o/0, o/1, o/2}{<3->}
  \acc{3}{b/1, b/2, b/3, o/0}{<4->}
  % left labels
  \begin{scope}[every node/.style={anchor=east, font=\footnotesize}]
    \node at (-0.65,0) {start: $X^{-\bar b}v=X^{-6}v$};
    \node[visible on=<2->] at (-0.65,-1*\dy) {$X^{-3}v$};
    \node[visible on=<3->] at (-0.65,-2*\dy) {$X^{-3}v$};
    \node[visible on=<4->] at (-0.65,-3*\dy) {$X^{-1}v$};
  \end{scope}
  % steps
  \begin{scope}[every node/.style={anchor=west, font=\footnotesize, align=left}]
    \draw[->, cNavy!70, line width=0.8pt, visible on=<2->] (1.5*\cw,-0.32) -- (1.5*\cw,-1*\dy+0.32);
    \node[visible on=<2->] at (3.75*\cw,-0.5*\dy)
      {$\BK_1$: $\bar a_1=3$, ${\color{cOrange!90!black}s_1=1}$ $\Rightarrow\times X^{3}$};
    \draw[->, cNavy!70, line width=0.8pt, visible on=<3->] (1.5*\cw,-1*\dy-0.32) -- (1.5*\cw,-2*\dy+0.32);
    \node[visible on=<3->] at (3.75*\cw,-1.5*\dy)
      {$\BK_2$: $\bar a_2=5$, ${\color{cBlue!90!black}s_2=0}$ $\Rightarrow$ unchanged};
    \draw[->, cNavy!70, line width=0.8pt, visible on=<4->] (1.5*\cw,-2*\dy-0.32) -- (1.5*\cw,-3*\dy+0.32);
    \node[visible on=<4->] at (3.75*\cw,-2.5*\dy)
      {$\BK_3$: $\bar a_3=2$, ${\color{cOrange!90!black}s_3=1}$ $\Rightarrow\times X^{2}$};
  \end{scope}
  \node[visible on=<5->, anchor=west, font=\footnotesize] at (3.75*\cw,-3*\dy)
    {output ${\color{cGreen!70!black}v_1}$, \ as $\bar\varphi=6-3-2=1$};
\end{tikzpicture}

\caption{The blind rotation of Example~\ref{ex:blindrot} ($n=3$, $N=4$,
$s=(1,0,1)$, $\bar b=6$, $\bar{\mathbf a}=(3,5,2)$) for any test polynomial
$v=v_0+v_1X+v_2X^2+v_3X^3$. Each row is the message of the accumulator (the
coefficients of the polynomial it encrypts) after one CMux. The server does
the same three CMux operations whatever $s$ is; only the key decides which of
them rotate. The final constant coefficient is $v_{\bar\varphi}=v_1$. With the
gate test polynomial ($v_j=\mu$), the rows are the ones in the table of the
example.}
\label{fig:blindrot}
\end{figure}


\subsection{Step 4: Sample Extraction}
\label{sec:extract}

Blind rotation gives a TRLWE ciphertext, which encrypts a whole table, but we
only want its constant coefficient, and the next gate expects an LWE
ciphertext. Luckily, one coefficient of a TRLWE ciphertext can be read as an
LWE ciphertext without any computation. Write the accumulator as $(a,b)$. The
constant coefficient of its phase $b-a\cdot s'$ is
\[
  b_0-a_0s'_0-\sum_{j=1}^{N-1}a_{N-j}s'_j,
\]
because, in the product $a\cdot s'$, the pairs of degrees that add up to $N$
end up in the constant term with their sign changed ($X^N=-1$). This is the
phase of the LWE ciphertext $\big((a_0,-a_{N-1},\dots,-a_1),\,b_0\big)$ under the
key formed by the coefficients of $s'$, seen as a vector of $N$ bits
(Figure~\ref{fig:extract}). Extraction only reorders coefficients and changes
some signs, so it adds no error.

% Sample extraction for N = 4.
\begin{figure}[!htb]
\centering
% Sample extraction for N = 4: the constant coefficient of a * s' is the inner
% product of s' with a rearranged and partly negated copy of a.
% Shared by the report and the slides.
\begin{tikzpicture}[font=\footnotesize, >=Stealth]
  \def\cw{1.05}
  \def\ch{0.5}
  \foreach \v [count=\j from 0] in {a_0,a_1,a_2,a_3}
    \node[draw=cNavy!40, fill=cGray!8, minimum width=\cw cm, minimum height=\ch cm] (a\j) at (\j*\cw,0) {$\v$};
  \foreach \v/\f [count=\j from 0] in {a_0/cBlue!14,-a_3/cOrange!16,-a_2/cOrange!16,-a_1/cOrange!16}
    \node[draw=cNavy!40, fill=\f, minimum width=\cw cm, minimum height=\ch cm] (e\j) at (\j*\cw,-1.6) {$\v$};
  \foreach \v [count=\j from 0] in {s'_0,s'_1,s'_2,s'_3}
    \node[draw=cNavy!40, fill=cGreen!12, minimum width=\cw cm, minimum height=\ch cm] (s\j) at (\j*\cw,-2.45) {$\v$};
  \draw[->, cBlue, line width=0.8pt] (a0.south) -- (e0.north);
  \draw[->, cOrange, line width=0.8pt] (a3.south) -- (e1.north);
  \draw[->, cOrange, line width=0.8pt] (a2.south) -- (e2.north);
  \draw[->, cOrange, line width=0.8pt] (a1.south) -- (e3.north);
  \node[anchor=east] at (-0.7,0) {TRLWE mask $a$};
  \node[anchor=east] at (-0.7,-1.6) {LWE mask $a'$};
  \node[anchor=east] at (-0.7,-2.45) {LWE key $=$ coefficients of $s'$};
  \foreach \j in {0,...,3} \node at (\j*\cw,-2.02) {$\times$};
  \node[anchor=west, align=left] at (4*\cw-0.2,-0.8)
    {$a\cdot s'$ has constant coefficient\\
     $a_0s'_0\;{\color{cOrange!90!black}-\;a_3s'_1-a_2s'_2-a_1s'_3}$\\[2pt]
     \scriptsize\color{cGray} (degrees adding up to $4=N$ wrap\\[-1pt]
     \scriptsize\color{cGray} around with a sign flip: $X^4=-1$)};
\end{tikzpicture}

\caption{Sample extraction for $N=4$. In the constant coefficient of
$a\cdot s'$, each $s'_j$ is multiplied by $a_{N-j}$, with a minus sign for
$j\ge1$. Reordering the mask and changing these signs gives an LWE ciphertext
$(a',b_0)$ whose phase, under the coefficients of $s'$, is the constant
coefficient of the TRLWE phase. Nothing is computed and no noise is added.}
\label{fig:extract}
\end{figure}


\subsection{Step 5: Key Switching}
\label{sec:ks}

\paragraph{Why it is needed.}
The extracted ciphertext has the right message but the wrong format: it is
encrypted under $s'$ and has size $N=1024$, while the inputs of every gate are
encrypted under $s$ and have size $n=630$. Without a conversion, the result
could not be used by another gate, because the bootstrapping key encrypts the
bits of $s$, not the bits of $s'$.

\paragraph{The idea.}
We want to keep the phase $b-\sum_{i=1}^{N}a_is'_i$. If the server had LWE
encryptions under $s$ of the products $a_is'_i$, it could subtract all of them
from the trivial ciphertext $(\mathbf0,b)$ and get an encryption under $s$ with
the same phase. It cannot store an encryption for every possible value of
$a_i$, but, as in the gadget decomposition, it can work digit by digit. We
write each $a_i$ with $t$ digits in base $\Bks=2^{\mathrm{basebit}}$,
$a_i\approx\sum_{j=1}^{t}a_{i,j}\Bks^{-j}$ with $a_{i,j}\in\{0,\dots,\Bks-1\}$. The
\emph{key-switching key} contains, for every coefficient $i$, level $j$ and digit
value $v$, the LWE ciphertext
$\KSK_{i,j,v}=\mathrm{LWE}_s\big(v\,s'_i/\Bks^{\,j}\big)$. Then
\[
  c'=(\mathbf{0},b)-\sum_{i=1}^{N}\sum_{j=1}^{t}\KSK_{i,j,a_{i,j}}
\]
has phase $b-\sum_is'_i\sum_ja_{i,j}\Bks^{-j}\approx b-\sum_ia_is'_i$ under $s$,
plus the errors of the ciphertexts we subtracted. Since the table has an entry
for every digit value, key switching only picks entries and subtracts them,
with no multiplication at all.

\begin{example}[One coefficient]
With $\Bks=4$ and $t=8$, the coefficient $a_i=0.7$ has base-4 digits
$(2,3,0,3,0,3,0,3)$, because $0.7=0.2303\ldots$ in base $4$. Its part of the key
switching costs five subtractions, one for each nonzero digit
(Figure~\ref{fig:kstable}). Doing the same for all $N=1024$ coefficients
completes the key switching.
\end{example}

% One slice of the key-switching table and the entries selected by a_i = 0.7.
\begin{figure}[!htb]
\centering
% Key switching: one slice (fixed coefficient i) of the key-switching table and
% the entries selected by the digits of a_i = 0.7. Shared by report and slides.
\begin{tikzpicture}[font=\footnotesize, >=Stealth]
  \def\cw{1.15}
  \def\ch{0.4}
  \foreach \v in {0,...,3} \node[font=\scriptsize, cGray] at (\v*\cw,0.42) {$v=\v$};
  \foreach \d [count=\j] in {2,3,0,3,0,3,0,3} {
    \node[font=\scriptsize, anchor=east] at (-0.75,-\j*\ch+\ch) {level $j=\j$};
    \node[font=\scriptsize, anchor=east, cBlue!85!black] at (-2.35,-\j*\ch+\ch) {$a_{i,\j}=\d$};
    \foreach \v in {0,...,3} {
      \ifnum\v=\d
        \ifnum\d=0
          \node[draw=cNavy!35, fill=cGray!25, minimum width=\cw cm, minimum height=\ch cm,
                inner sep=0pt, font=\scriptsize] at (\v*\cw,-\j*\ch+\ch) {skip};
        \else
          \node[draw=cNavy!60, fill=cGreen!35, minimum width=\cw cm, minimum height=\ch cm,
                inner sep=0pt, font=\scriptsize, line width=0.8pt] at (\v*\cw,-\j*\ch+\ch) {subtract};
        \fi
      \else
        \node[draw=cNavy!25, fill=white, minimum width=\cw cm, minimum height=\ch cm,
              inner sep=0pt] at (\v*\cw,-\j*\ch+\ch) {};
      \fi
    }
  }
  \node[anchor=west, align=left] at (3.75*\cw,-0.9)
    {entry $(i,j,v)$:\\$\mathrm{LWE}_s\big(v\,s'_i/4^{j}\big)$};
  \node[anchor=west, align=left] at (3.75*\cw,-2.2)
    {$a_i=0.7=\texttt{0xB3333333}$\\in base 4: $0.23030303\ldots$\\[2pt]
     $\Rightarrow$ 5 subtractions, no products};
\end{tikzpicture}

\caption{Key switching for one coefficient $a_i$ of the extracted ciphertext.
For each of the $t=8$ digit positions $j$, the table has an LWE encryption
under $s$ of $v\cdot s'_i/4^j$ for every possible digit $v$. The base-4 digits of
$a_i$ choose one entry per position; the chosen entries are subtracted from
$(\mathbf 0,b)$, and zero digits are skipped. Doing this for all $N=1024$
coefficients replaces the part $\sum_ia_is'_i$ of the phase by the same value
computed under the key $s$.}
\label{fig:kstable}
\end{figure}


Algorithm~\ref{alg:bootstrap} puts the five steps together.

\begin{algorithm}[!htb]
\caption{Gate bootstrapping of an LWE ciphertext $c=(\mathbf{a},b)$ under $s$}
\label{alg:bootstrap}
\begin{algorithmic}[1]
\Require $\mu\in\T$; $\BK_i=\mathrm{TRGSW}_{s'}(s_i)$; $\KSK$
\Ensure LWE ciphertext under $s$ of $+\mu$ if $\varphi_s(c)\in(0,\frac12)$, of $-\mu$ otherwise
\State $\bar b\gets\round{2Nb}$;\quad $v\gets\mu\,(1+X+\dots+X^{N-1})$
  \Comment{steps 1 and 2}
\State $\mathrm{ACC}\gets(0,\;X^{-\bar b}\cdot v)$ \Comment{trivial TRLWE ciphertext}
\For{$i=1,\dots,n$} \Comment{step 3: blind rotation}
  \State $\bar a_i\gets\round{2Na_i}$
  \If{$\bar a_i\neq0$}
    \State $\mathrm{ACC}\gets\mathrm{ACC}+\BK_i\ext\big((X^{\bar a_i}-1)\cdot\mathrm{ACC}\big)$
      \Comment{CMux, Eq.~\eqref{eq:cmux}}
  \EndIf
\EndFor
\State $(a,b)\gets\mathrm{ACC}$;\quad
       $c'\gets\big((a_0,-a_{N-1},\dots,-a_1),\,b_0\big)$
       \Comment{step 4: key $s'$}
\State \Return $\mathrm{KeySwitch}(c',\KSK)$ \Comment{step 5: key $s$, size $n$}
\end{algorithmic}
\end{algorithm}

\subsection{Why the Noise Is Reset}
\label{sec:bsnoise}

Where does the error of the output come from? Blind rotation starts from a
trivial ciphertext with \emph{no} noise, so all the error it collects comes
from the bootstrapping key, through the $n$ external products. Key switching
then adds the errors of the table entries it subtracts. The error of the input
does not enter any of these operations: it only affects which slice
$\bar\varphi$ is chosen in step 1. So, as long as the input error plus the
modulus-switching error keeps the phase in the correct half of the circle, the
output is a correct encryption, and its error depends only on the evaluation
key. Section~\ref{sec:noise} shows this with a deliberately noisy input.

To put numbers on it, each of the $n$ CMux steps adds the error of one external
product (Proposition~\ref{prop:extprod}), and key switching adds the error of up
to $tN$ table entries plus the rounding of each $a_i$ to $t$ digits:
\begin{align}
  \vartheta_{\mathrm{BR}}&\le n\big(2\ell N\beta^2\vartheta_{\BK}+(1+N)\varepsilon^2\big),
  \label{eq:varbr}\\
  \vartheta_{\mathrm{KS}}&\le tN\vartheta_{\KSK}+N\cdot2^{-2(t\cdot\mathrm{basebit}+1)}.
  \label{eq:varks}
\end{align}
With the parameters of Table~\ref{tab:params}, $\vartheta_{\mathrm{BR}}\le
\BoundBR$ and $\vartheta_{\mathrm{BR}}+\vartheta_{\mathrm{KS}}\le\BoundTotal$,
which is a standard deviation below $0.005$. These bounds have the same form as
the analysis of \cite{chillotti2020tfhe}; Section~\ref{sec:noise} compares them
with our measurements.

\subsection{From a Sign Function to Logic Gates}
\label{sec:gates}

With the gate test polynomial, the bootstrapping returns $+\frac18$
(\textsf{true}) when the phase of its input is in the upper half of the circle,
and $-\frac18$ (\textsf{false}) when it is in the lower half. To compute a gate
with two inputs, it is therefore enough to build, from the two inputs, a
ciphertext whose phase is in the upper half exactly when the output of the
gate should be \textsf{true}. Linear combinations are free
(Section~\ref{sec:limits}), so we only need to choose the right one.

\paragraph{NAND.}
Take $(\mathbf{0},\frac18)-c_x-c_y$, the trivial ciphertext of $\frac18$ minus
the two inputs. Since the inputs are $\pm\frac18$, its phase is
\[
  \tfrac18-(\pm\tfrac18)-(\pm\tfrac18)\in
  \big\{\underbrace{\tfrac38}_{(0,0)},\;
        \underbrace{\tfrac18}_{(0,1),(1,0)},\;
        \underbrace{-\tfrac18}_{(1,1)}\big\},
\]
which is in the upper half for every input except $(1,1)$. This is the truth
table of NAND (Figure~\ref{fig:gates}a). All these points are at distance
$\frac18$ from the closest boundary ($0$ or $\frac12$). This distance is the
\emph{margin}: the input errors plus the modulus-switching error must stay
below it for the gate to give the right answer.

\begin{example}[One NAND gate]
Let $c_x$ and $c_y$ both encrypt \textsf{true}, with phases $\frac18+e_x$ and
$\frac18+e_y$. The combination $(\mathbf0,\frac18)-c_x-c_y$ has phase
$\frac18-\frac18-\frac18-(e_x+e_y)=-\frac18-(e_x+e_y)$, which is in the lower half
as long as $|e_x+e_y|$ plus the modulus-switching error is less than $\frac18$.
So the bootstrapping returns an encryption of $-\frac18$, that is,
\textsf{false}, with fresh noise that no longer depends on $e_x$ and $e_y$.
\end{example}

\paragraph{Other gates.}
The same reasoning gives AND, OR and NOR, by changing the constant and the
signs (Table~\ref{tab:gates}). XOR needs one more idea. The phase of $c_x+c_y$
is $-\frac14$, $0$ or $+\frac14$, and the middle value, which corresponds to
different inputs, is exactly on the boundary. Doubling fixes this: $2(c_x+c_y)$
has phase $-\frac12$, $0$ or $+\frac12$, and since $-\frac12$ and $+\frac12$ are
the same point of the torus, the two cases with equal inputs now fall on the
same point. Adding $\frac14$ moves them to $-\frac14$ and the case with different
inputs to $+\frac14$ (Figure~\ref{fig:gates}b). Doubling also doubles the input
errors, but the margin doubles too. NOT does not need a bootstrapping at all:
$-c$ swaps $\pm\frac18$ and only changes the sign of the error.

\begin{table}[!htb]
\centering
\caption{Linear combination computed before the bootstrapping of each gate.}
\label{tab:gates}
\small
\renewcommand{\arraystretch}{1.35}
\begin{tabular}{@{}lll@{\qquad}lll@{}}
\toprule
Gate & Combination & Margin & Gate & Combination & Margin\\
\midrule
NAND & $(\mathbf{0},+\frac18)-c_x-c_y$ & $\frac18$ &
AND  & $(\mathbf{0},-\frac18)+c_x+c_y$ & $\frac18$\\[2pt]
OR   & $(\mathbf{0},+\frac18)+c_x+c_y$ & $\frac18$ &
NOR  & $(\mathbf{0},-\frac18)-c_x-c_y$ & $\frac18$\\[2pt]
XOR  & $(\mathbf{0},+\frac14)+2(c_x+c_y)$ & $\frac14$ &
XNOR & $(\mathbf{0},-\frac14)-2(c_x+c_y)$ & $\frac14$\\
\bottomrule
\end{tabular}
\end{table}

% Phases at the input of the bootstrapping for NAND and XOR.
\begin{figure}[!htb]
\centering
\begin{tikzpicture}[font=\small]
  \def\R{1.65}
  % ---------------- (a) NAND ----------------
  \begin{scope}
    \fill[cBlue!14] (0,0) -- (\R,0) arc (0:180:\R) -- cycle;
    \fill[cOrange!16] (0,0) -- (-\R,0) arc (180:360:\R) -- cycle;
    \draw[line width=0.9pt] (0,0) circle (\R);
    \draw[cGray, densely dotted] (-\R-0.1,0) -- (\R+0.1,0);
    \fill[cBlue] (135:\R) circle (2.3pt);
    \fill[cBlue] (45:\R) circle (2.3pt);
    \fill[cOrange] (-45:\R) circle (2.3pt);
    \node[anchor=south east, font=\footnotesize, align=right] at (135:\R+0.02)
      {$+\tfrac38$: $(0,0)$};
    \node[anchor=south west, font=\footnotesize, align=left] at (45:\R+0.02)
      {$+\tfrac18$: $(0,1),(1,0)$};
    \node[anchor=north west, font=\footnotesize] at (-45:\R+0.02) {$-\tfrac18$: $(1,1)$};
    \node[font=\footnotesize, cBlue!80!black] at (0,0.55) {output $+\tfrac18$};
    \node[font=\footnotesize, cOrange!80!black] at (0,-0.55) {output $-\tfrac18$};
    \node[right, font=\scriptsize] at (\R+0.1,0) {$0$};
    \node[left, font=\scriptsize] at (-\R-0.1,0) {$\tfrac12$};
    \node[font=\footnotesize] at (0,-\R-1.0) {$\varphi\big((0,\tfrac18)-c_x-c_y\big)$};
    \node at (-\R-1.2, 0.9) {(a)};
  \end{scope}
  % ---------------- (b) XOR ----------------
  \begin{scope}[xshift=7.6cm]
    \fill[cBlue!14] (0,0) -- (\R,0) arc (0:180:\R) -- cycle;
    \fill[cOrange!16] (0,0) -- (-\R,0) arc (180:360:\R) -- cycle;
    \draw[line width=0.9pt] (0,0) circle (\R);
    \draw[cGray, densely dotted] (-\R-0.1,0) -- (\R+0.1,0);
    \fill[cBlue] (90:\R) circle (2.3pt);
    \fill[cOrange] (-90:\R) circle (2.3pt);
    \node[above, font=\footnotesize] at (90:\R+0.02) {$+\tfrac14$: $(0,1),(1,0)$};
    \node[below, font=\footnotesize, align=center] at (-90:\R+0.02)
      {$\tfrac12+\tfrac14=-\tfrac14$: $(0,0),(1,1)$};
    \node[font=\footnotesize, cBlue!80!black] at (0,0.55) {output $+\tfrac18$};
    \node[font=\footnotesize, cOrange!80!black] at (0,-0.55) {output $-\tfrac18$};
    \node[right, font=\scriptsize] at (\R+0.1,0) {$0$};
    \node[left, font=\scriptsize] at (-\R-0.1,0) {$\tfrac12$};
    \node[font=\footnotesize] at (0,-\R-1.0) {$\varphi\big((0,\tfrac14)+2(c_x+c_y)\big)$};
    \node at (-\R-1.2, 0.9) {(b)};
  \end{scope}
\end{tikzpicture}
\caption{Phases at the input of the bootstrapping, without noise. (a) NAND:
every point is at distance $\frac18$ from the closest boundary ($0$ or
$\frac12$); this distance is the noise margin. (b) XOR: with $c_x+c_y$, different
inputs would fall on the boundary. After doubling, the two cases with equal
inputs fall on the same point $\pm\frac12$, and adding $\frac14$ separates the
two groups with a margin of $\frac14$.}
\label{fig:gates}
\end{figure}


\paragraph{MUX.}
The multiplexer $\mathrm{MUX}(s,x,y)$, ``$x$ if $s$ is true, otherwise $y$'', can
be written as $(s\wedge x)\vee(\neg s\wedge y)$, which would need three
bootstrappings. But the two ANDs are never true at the same time, so their OR
can be replaced by a simple sum. We compute both ANDs with bootstrappings
\emph{without} key switching (so their outputs are under $s'$), add the two
results, which gives $-\frac14$ or $0$, add $\frac18$ to get $-\frac18$ or
$+\frac18$, and do a single key switching at the end (Table~\ref{tab:mux}). A
MUX therefore costs two blind rotations and one key switching.

\begin{table}[!htb]
\centering
\caption{The MUX gate on all inputs. $u_1$ and $u_2$ are the results of the two
bootstrappings without key switching. At most one of them is $+\frac18$, so
$u_1+u_2+\frac18$ is always $\pm\frac18$.}
\label{tab:mux}
\small
\renewcommand{\arraystretch}{1.35}
\begin{tabular}{@{}ccc@{\qquad}cc@{\qquad}cc@{}}
\toprule
$s$ & $x$ & $y$ & $u_1=\mathrm{BR}(s\wedge x)$ & $u_2=\mathrm{BR}(\neg s\wedge y)$
  & $u_1+u_2+\frac18$ & output\\
\midrule
0 & 0 & 0 & $-\frac18$ & $-\frac18$ & $-\frac18$ & $0=y$\\
0 & 0 & 1 & $-\frac18$ & $+\frac18$ & $+\frac18$ & $1=y$\\
0 & 1 & 0 & $-\frac18$ & $-\frac18$ & $-\frac18$ & $0=y$\\
0 & 1 & 1 & $-\frac18$ & $+\frac18$ & $+\frac18$ & $1=y$\\
1 & 0 & 0 & $-\frac18$ & $-\frac18$ & $-\frac18$ & $0=x$\\
1 & 0 & 1 & $-\frac18$ & $-\frac18$ & $-\frac18$ & $0=x$\\
1 & 1 & 0 & $+\frac18$ & $-\frac18$ & $+\frac18$ & $1=x$\\
1 & 1 & 1 & $+\frac18$ & $-\frac18$ & $+\frac18$ & $1=x$\\
\bottomrule
\end{tabular}
\end{table}

\subsection{An Encrypted Adder}
\label{sec:adder}

Every gate gives back a ciphertext in the same format as its inputs and with
fresh noise, so gates can be connected into circuits of any size and depth. As
a complete example, we built a circuit that adds two encrypted integers.

\paragraph{Encrypting an integer.}
TFHE encrypts bits, so an $n$-bit integer is encrypted bit by bit, as $n$
separate LWE ciphertexts, from the least significant bit to the most
significant one. For example, $x=11=1011_2$ is stored as the four ciphertexts
$(\mathrm{Enc}(1),\mathrm{Enc}(1),\mathrm{Enc}(0),\mathrm{Enc}(1))$.

\paragraph{Adding by hand in binary.}
The circuit follows the method we learn in school: add the two numbers column
by column, from right to left, and carry a $1$ to the next column when the
column sum is $2$ or more. For $11+6$:
\begin{center}
\small
\begin{tabular}{@{}r@{\qquad}ccccc@{}}
carries & 1 & 1 & 1 & 0 & \\
$x=11$ &   & 1 & 0 & 1 & 1\\
$y=6$  &   & 0 & 1 & 1 & 0\\
\midrule
$x+y=17$ & 1 & 0 & 0 & 0 & 1\\
\end{tabular}
\end{center}
In each column $i$ there are three input bits, $x_i$, $y_i$ and the incoming
carry $c_i$, and two output bits, the sum bit $\mathit{sum}_i$ and the outgoing
carry $c_{i+1}$. A circuit that computes one column is called a \emph{full
adder}, and Table~\ref{tab:fulladder} gives its truth table.

\begin{table}[!htb]
\centering
\caption{Truth table of a full adder, with the formulas used by our circuit
($p=x\oplus y$).}
\label{tab:fulladder}
\small
\begin{tabular}{@{}ccc@{\qquad}cc@{\qquad}c@{}}
\toprule
$x_i$ & $y_i$ & $c_i$ & $\mathit{sum}_i$ & $c_{i+1}$ & which case\\
\midrule
0 & 0 & 0 & 0 & 0 & $x_i=y_i$: carry is $x_i=0$\\
0 & 0 & 1 & 1 & 0 & $x_i=y_i$: carry is $x_i=0$\\
0 & 1 & 0 & 1 & 0 & $x_i\ne y_i$: carry is $c_i=0$\\
0 & 1 & 1 & 0 & 1 & $x_i\ne y_i$: carry is $c_i=1$\\
1 & 0 & 0 & 1 & 0 & $x_i\ne y_i$: carry is $c_i=0$\\
1 & 0 & 1 & 0 & 1 & $x_i\ne y_i$: carry is $c_i=1$\\
1 & 1 & 0 & 0 & 1 & $x_i=y_i$: carry is $x_i=1$\\
1 & 1 & 1 & 1 & 1 & $x_i=y_i$: carry is $x_i=1$\\
\bottomrule
\end{tabular}
\end{table}

\paragraph{From the truth table to gates.}
The sum bit is $1$ when an odd number of the three inputs are $1$, so
$\mathit{sum}_i=x_i\oplus y_i\oplus c_i$. For the carry, the last column of
Table~\ref{tab:fulladder} shows a simple rule. When $x_i=y_i$, the outgoing
carry is equal to $x_i$, whatever the incoming carry is: two $1$s always produce
a carry and two $0$s never do. When $x_i\ne y_i$, the column sum is $1+c_i$, so
the outgoing carry is just the incoming one. This is a choice controlled by
$p_i=x_i\oplus y_i$, so
\[
  \mathit{sum}_i=p_i\oplus c_i,\qquad c_{i+1}=\mathrm{MUX}(p_i;\,c_i,\,x_i).
\]
The value $p_i$ is used twice, so one column needs only two XOR gates and one
MUX (Figure~\ref{fig:adder}), that is, four blind rotations. Adding two $n$-bit
numbers takes $n$ columns and $4n$ blind rotations, and the result has $n+1$
bits, the last one being the final carry. The first carry, $c_0=0$, is a trivial
ciphertext of \textsf{false}, which the server can create without any key.

% One cell of the homomorphic ripple-carry adder.
\begin{figure}[!htb]
\centering
\begin{tikzpicture}[font=\footnotesize, >=stealth,
    gate/.style={draw, rounded corners=2pt, fill=cBlue!10, minimum width=1.2cm,
                 font=\footnotesize\sffamily, inner sep=0pt},
    wire/.style={->, line width=0.6pt},
    bridge/.style={preaction={draw, white, line width=3.2pt}},
    dot/.style={circle, fill, inner sep=1.1pt}]
  % gates
  \node[gate, minimum height=0.9cm] (x1) at (2.1, 0.7) {XOR};
  \node[gate, minimum height=0.9cm] (x2) at (5.2, 0.55) {XOR};
  \node[gate, minimum height=1.2cm, fill=cOrange!14] (mx) at (5.2,-1.35) {MUX};
  % inputs
  \node[anchor=east] at (0, 1.0)   {$y_i$};
  \node[anchor=east] at (0, 0.4)   {$x_i$};
  \node[anchor=east] at (0,-1.05)  {$c_i$};
  \draw[wire] (0, 1.0) -- (1.5, 1.0);
  \draw[wire] (0, 0.4) -- (1.5, 0.4);
  % c_i: to the MUX (if 1) and up to the second XOR
  \draw[wire] (0,-1.05) -- (4.6,-1.05);
  \node[dot] at (3.9,-1.05) {};
  \draw[wire] (3.9,-1.05) -- (3.9, 0.3) -- (4.6, 0.3);
  % x_i: branch down to the MUX (if 0), bridging over c_i
  \node[dot] at (0.75, 0.4) {};
  \draw[wire, bridge] (0.75, 0.4) -- (0.75,-1.65) -- (4.6,-1.65);
  % p_i = x_i XOR y_i: to the second XOR and to the selector
  \draw[wire] (2.7, 0.7) -- node[above, pos=0.3] {$p_i$} (4.6, 0.7);
  \node[dot] at (3.3, 0.7) {};
  \draw[wire, bridge] (3.3, 0.7) -- (3.3,-0.4) -- (5.2,-0.4) -- (5.2,-0.75);
  % MUX input labels
  \node[font=\scriptsize, anchor=west] at (4.62,-1.05) {1};
  \node[font=\scriptsize, anchor=west] at (4.62,-1.65) {0};
  \node[font=\scriptsize, anchor=south west] at (5.2,-0.62) {sel};
  % outputs
  \draw[wire] (x2.east) -- (6.7, 0.55) node[right] {$\mathit{sum}_i$};
  \draw[wire] (mx.east) -- (6.7,-1.35) node[right] {$c_{i+1}$};
\end{tikzpicture}
\caption{One column of the adder (a full adder): $\mathit{sum}_i=p_i\oplus c_i$
and $c_{i+1}=\mathrm{MUX}(p_i;\,c_i,\,x_i)$, with $p_i=x_i\oplus y_i$. When
$x_i=y_i$ the outgoing carry is $x_i$, whatever $c_i$ is; otherwise the incoming
carry $c_i$ goes through. Each column uses two XOR gates and one MUX, that is,
four blind rotations.}
\label{fig:adder}
\end{figure}


\begin{example}[$11+6$ on encrypted bits]
With $x=11$, $y=6$ and $c_0=0$, the circuit goes through the four columns as
follows (bits from the least significant one):
\begin{center}
\small
\begin{tabular}{@{}ccccccc@{}}
\toprule
column $i$ & $x_i$ & $y_i$ & $c_i$ & $p_i=x_i\oplus y_i$ & $\mathit{sum}_i=p_i\oplus c_i$
  & $c_{i+1}=\mathrm{MUX}(p_i;c_i,x_i)$\\
\midrule
0 & 1 & 0 & 0 & 1 & 1 & $c_0=0$\\
1 & 1 & 1 & 0 & 0 & 0 & $x_1=1$\\
2 & 0 & 1 & 1 & 1 & 0 & $c_2=1$\\
3 & 1 & 0 & 1 & 1 & 0 & $c_3=1$\\
\bottomrule
\end{tabular}
\end{center}
The output bits are $(\mathit{sum}_0,\dots,\mathit{sum}_3,c_4)=(1,0,0,0,1)$,
i.e., $10001_2=17$, the same result as by hand. On encrypted data, every entry of
this table is a ciphertext, and the server never sees any of these bits.
Because of this, the server cannot stop early or skip a column when there is
no carry: it always evaluates the same $12$ gates, whatever the inputs are.
\end{example}

%% ===========================================================================
\section{Implementation}
\label{sec:impl}

\subsection{Organization}

The library is made of eleven C++ header files with \code{inline} functions, so
each test or program is a single source file and nothing has to be linked.
The headers form a chain in which each layer includes the previous ones
(Figure~\ref{fig:arch}). Each header has its own unit test, which checks it
against a simple, independent reference before the next layer is built on top
of it (Table~\ref{tab:tests}). The library has \LibLines{} lines of code, not
counting comments and blank lines. We use the parameters of the 128-bit
gate-bootstrapping set of the original library (Table~\ref{tab:params}), so
that the comparison in Section~\ref{sec:eval} is fair.

% Module structure of the implementation: each header includes the ones below it.
\begin{figure}[!htb]
\centering
\begin{tikzpicture}[font=\footnotesize, >=stealth,
    mod/.style={draw, rounded corners=2pt, fill=#1, minimum height=0.5cm,
                minimum width=2.4cm, inner sep=2pt, font=\footnotesize\ttfamily},
    mod/.default=cBlue!10,
    app/.style={mod=cGreen!14, minimum width=2.0cm},
    dep/.style={->, cGray!85!black, line width=0.6pt},
    ann/.style={font=\scriptsize, cGray!90!black, align=left}]
  \def\s{0.82}
  \node[mod=cGray!10] (torus) at (0,0)        {torus.hpp};
  \node[mod] (lwe)   at (-2.9,2.5*\s)          {lwe.hpp};
  \node[mod] (poly)  at ( 2.8,1*\s)            {poly.hpp};
  \node[mod] (fft)   at ( 2.8,2*\s)            {fft.hpp};
  \node[mod] (trlwe) at ( 2.8,3*\s)            {trlwe.hpp};
  \node[mod] (trgsw) at ( 2.8,4*\s)            {trgsw.hpp};
  \node[mod=cGray!10] (params) at (-2.9,5*\s)  {params.hpp};
  \node[mod] (keys)  at (0,5*\s)               {keys.hpp};
  \node[mod] (br)    at (0,6*\s)               {blindrotate.hpp};
  \node[mod] (boot)  at (0,7*\s)               {bootstrap.hpp};
  \node[mod] (gates) at (0,8*\s)               {gates.hpp};
  \node[app] (adder) at (-1.3,9.2*\s)          {adder.cpp};
  \node[app] (bench) at ( 1.3,9.2*\s)          {bench.cpp};
  \draw[dep] (torus.west) -| (lwe.south);
  \draw[dep] (torus.east) -| (poly.south);
  \draw[dep] (poly) -- (fft);
  \draw[dep] (fft) -- (trlwe);
  \draw[dep] (trlwe) -- (trgsw);
  \draw[dep] (lwe.north) |- (-0.6,4.5*\s) -- ([xshift=-0.6cm]keys.south);
  \draw[dep] (params) -- (keys);
  \draw[dep] (trgsw.north) |- (0.6,4.5*\s) -- ([xshift=0.6cm]keys.south);
  \draw[dep] (keys) -- (br);
  \draw[dep] (br) -- (boot);
  \draw[dep] (boot) -- (gates);
  \draw[dep] (gates) -- (adder);
  \draw[dep] (gates) -- (bench);
  \node[ann, anchor=west] at (4.15,0)       {$\mathbb{T}$ as \texttt{uint32\_t}, encoding, sampling};
  \node[ann, anchor=west] at (4.15,1*\s)    {$\mathbb{Z}_N[X]$, $\mathbb{T}_N[X]$, product by $X^e$};
  \node[ann, anchor=west] at (4.15,2*\s)    {negacyclic FFT};
  \node[ann, anchor=west] at (4.15,3*\s)    {TRLWE samples};
  \node[ann, anchor=west] at (4.15,4*\s)    {gadget, TRGSW, $\boxdot$, CMux};
  \node[ann, right=6pt of br]    {modulus switching, blind rotation, extraction};
  \node[ann, right=6pt of boot]  {key switching, bootstrapping, bit encoding};
  \node[ann, right=6pt of gates] {NAND, AND, OR, NOR, XOR, XNOR, NOT, MUX};
  \node[ann, anchor=east] at (-4.2,2.5*\s) {LWE samples};
  \node[ann, anchor=east] at (-4.2,5*\s)   {parameter set};
  \node[ann, right=6pt of keys] {secret and evaluation keys};
\end{tikzpicture}
\caption{Structure of the library. Each box is a header file in \texttt{src/}
with its own unit test in \texttt{tests/}; an arrow goes from a header to the
headers that include it. The programs in \texttt{apps/} only use
\texttt{gates.hpp}.}
\label{fig:arch}
\end{figure}


\begin{table}[!htb]
\centering
\caption{Parameters (the same as the 128-bit set of the original library).}
\label{tab:params}
\small
\begin{tabular}{@{}lll@{}}
\toprule
Symbol & Value & Role\\
\midrule
$n$ & $630$ & size of the LWE key; number of external products per bootstrapping\\
$N$ & $1024$ & degree of the polynomials; the phase is rounded to $2N$ positions\\
$\ell$, $\Bg$ & $3$, $2^7$ & number of digits and base of the gadget decomposition\\
$t$, $\Bks$ & $8$, $2^2$ & number of digits and base of the key switching\\
$\sigma_{\mathrm{LWE}}$ & $2^{-15}$ & noise of LWE ciphertexts and of $\KSK$\\
$\sigma_{\BK}$ & $2^{-25}$ & noise of the TRGSW ciphertexts of $\BK$\\
\bottomrule
\end{tabular}
\end{table}

\begin{table}[!htb]
\centering
\caption{Unit tests. \code{make test} runs them in order and stops at the first
failed check.}
\label{tab:tests}
\small
\begin{tabular}{@{}ll@{}}
\toprule
Test & What it checks\\
\midrule
\code{test01\_torus} & conversions, wrap-around, decoding in $\Z_8$, Gaussian noise\\
\code{test02\_lwe} & decryption, additions, noise of a sum growing by $\sqrt2$\\
\code{test03\_poly} & $X^N=-1$, $X^{2N}=1$, rotation against the direct product\\
\code{test04\_fft} & FFT and back; products identical to the direct ones\\
\code{test05\_trlwe} & TRLWE decryption, addition, rotation by $X^{100}$\\
\code{test06\_trgsw} & digit range and error, $\mathrm{TRGSW}(m)\ext d$, CMux\\
\code{test07\_keys} & key sizes; each $\BK_i$ selects correctly; $\KSK$ contents\\
\code{test08\_blindrotate} & rotation without noise, sign function, four-entry table\\
\code{test09\_bootstrap} & key switching, bootstrappings in a row, noise reduction\\
\code{test10\_gates} & truth tables, 50 NAND gates in a row, majority circuit\\
\bottomrule
\end{tabular}
\end{table}

\subsection{Torus Arithmetic}
\label{sec:torusimpl}

A computer cannot store real numbers exactly, so implementations use a
discretized torus: the 32-bit integer $t$ stands for the real number
$t/2^{32}$. With this choice the reduction modulo $1$ costs nothing, because
unsigned 32-bit arithmetic already wraps around modulo $2^{32}$.

\begin{example}[Torus arithmetic in 32 bits]
$0.75+0.5=1.25\equiv0.25$ becomes
$\texttt{0xC0000000}+\texttt{0x80000000}=\texttt{0x1\,4000\,0000}$, and keeping
the 32 low bits gives $\texttt{0x40000000}=2^{30}$, which is $0.25$. In the same
way, $3\times0.375\equiv0.125$ becomes
$3\times\texttt{0x60000000}=\texttt{0x1\,2000\,0000}\mapsto\texttt{0x20000000}$.
\end{example}

A torus element is therefore a \code{uint32\_t}, and three rules of C++ do all
the work. First, unsigned arithmetic is defined modulo $2^{32}$, which is
exactly what the torus needs. Signed overflow, on the other hand, is undefined
behavior, and an optimizing compiler may break code that relies on it. Second,
in an expression like \code{k * t} with a signed \code{k}, the integer is
converted to unsigned before the product, which still gives the right result
modulo $2^{32}$ when $k$ is negative. Third, \code{static\_cast<int32\_t>(t)}
reads the same bits as a signed integer, that is, as the representative of $t$
in $[-\frac12,\frac12)$. We use this to measure noise (a small negative error
reads as $-0.001$ and not as $0.999$) and to decrypt a bit (by the sign of its
phase). The bit $\mu=\frac18$ of the gates is the 32-bit value
\texttt{0x20000000}, so the gate test polynomial is built with the single line
\code{TorusPoly testv(N, MU)}.

\subsection{Fast Polynomial Products: the Negacyclic FFT}
\label{sec:fft}

\paragraph{Why we need it.}
Every external product multiplies $2\ell=6$ digit polynomials by the two
polynomials of a TRGSW row, so it computes $12$ products in $\ZN\times\TN$.
A gate does about $n=630$ external products, which makes about $7500$
polynomial products per gate. If we compute each product term by term, it
costs $N^2\approx10^6$ multiplications, and a single gate would take several
seconds. Instead, we use the well-known idea of multiplying polynomials through
their values, adapted to polynomials modulo $X^N+1$. We explain it in four
steps.

\paragraph{Step 1: multiply values instead of coefficients.}
A polynomial of degree less than $N$ is completely determined by its values at
$N$ different points, and the value of a product is the product of the values,
$(a\cdot b)(x)=a(x)\,b(x)$. So we can evaluate both factors at enough points,
multiply the values one by one, and then recover the coefficients of the
product from its values (interpolation).

\begin{example}[Multiplying through values]
Take $a(x)=1+2x$ and $b(x)=3+x$. Their product has degree $2$, so three points
are enough; we use $0$, $1$ and $-1$. The values are $a=(1,3,-1)$ and
$b=(3,4,2)$, and multiplying them one by one gives $(3,12,-2)$. The only
polynomial of degree $2$ with these values is $3+7x+2x^2$, which is indeed
$(1+2x)(3+x)$.
\end{example}

Multiplying the values costs only $N$ operations. Evaluating and interpolating
in a direct way would still cost about $N^2$ operations, but with the fast
Fourier transform (FFT) \cite{cooley1965fft} both cost about $N\log N$, if the
points are well chosen.

\paragraph{Step 2: choose points where the reduction is automatic.}
We do not want the full product $a\cdot b$, which has degree up to $2N-2$, but
its remainder modulo $X^N+1$. If we evaluate at a root $\omega$ of $X^N+1$, the
difference disappears: since $\omega^N=-1$, the product and its remainder have
the same value at $\omega$, because they differ by a multiple of $X^N+1$. The
$N$ roots of $X^N+1$ are the points $e^{i\pi(2k+1)/N}$ of the unit circle, the
odd multiples of $\pi/N$ (Figure~\ref{fig:roots}). Moreover, a polynomial with
real coefficients takes conjugate values at conjugate points, so we only need
one root of each conjugate pair. That leaves $N/2$ complex values, and we use
the roots $\omega_k=e^{i\pi(4k+1)/N}$ for $k=0,\dots,N/2-1$.

% Evaluation points of the negacyclic FFT for N = 8.
\begin{figure}[!htb]
\centering
% Roots of X^8 + 1 on the unit circle: the four evaluation points of the
% negacyclic FFT and their conjugates. Shared by the report and the slides.
\begin{tikzpicture}[font=\footnotesize, >=Stealth]
  \def\rad{1.65}
  \draw[->, cGray] (-2.3,0) -- (2.4,0) node[right] {Re};
  \draw[->, cGray] (0,-2.2) -- (0,2.3) node[above] {Im};
  \draw[cNavy!60, line width=0.8pt] (0,0) circle (\rad);
  \foreach \a/\b in {22.5/337.5, 112.5/247.5, 202.5/157.5, 292.5/67.5}
    \draw[cGray, densely dotted] (\a:\rad) -- (\b:\rad);
  \foreach \a/\k in {22.5/0, 112.5/1, 202.5/2, 292.5/3} {
    \fill[cBlue] (\a:\rad) circle (2.8pt);
    \node[cBlue!85!black] at (\a:\rad+0.42) {$\omega_{\k}$};
  }
  \foreach \a/\k in {337.5/0, 247.5/1, 157.5/2, 67.5/3} {
    \draw[cOrange, fill=white, line width=0.9pt] (\a:\rad) circle (2.6pt);
    \node[cOrange!85!black] at (\a:\rad+0.42) {$\bar\omega_{\k}$};
  }
  \draw[cGray] (0.6,0) arc (0:22.5:0.6);
  \node[cGray, font=\scriptsize] at (11:0.95) {$\frac{\pi}{8}$};
\end{tikzpicture}

\caption{The roots of $X^8+1$, at the odd multiples of $\pi/8$ on the unit
circle. The negacyclic FFT evaluates a polynomial at the four points
$\omega_k=e^{i\pi(4k+1)/N}$ (filled). The other four roots (hollow, joined by
dotted lines) are their complex conjugates, where a polynomial with real
coefficients takes conjugate values, so they do not need to be computed.}
\label{fig:roots}
\end{figure}


\begin{example}[A negacyclic product through values, $N=4$]
Let $a=1+X$ and $b=X^3$. Computed directly, $a\cdot b=X^3+X^4=-1+X^3$, so the
result is $r=[-1,0,0,1]$. The two points are $\omega_0=e^{i\pi/4}$ and
$\omega_1=e^{i5\pi/4}$:
\begin{center}
\small
\begin{tabular}{@{}lcc@{}}
\toprule
 & $\omega_0=e^{i\pi/4}$ & $\omega_1=e^{i5\pi/4}$\\
\midrule
$a(\omega)=1+\omega$ & $1.707+0.707i$ & $0.293-0.707i$\\
$b(\omega)=\omega^3$ & $-0.707+0.707i$ & $0.707-0.707i$\\
$a(\omega)\,b(\omega)$ & $-1.707+0.707i$ & $-0.293-0.707i$\\
$r(\omega)=-1+\omega^3$ & $-1.707+0.707i$ & $-0.293-0.707i$\\
\bottomrule
\end{tabular}
\end{center}
The products of the values are already the values of the reduced product $r$,
because $\omega^4=-1$ at both points.
\end{example}

\paragraph{Step 3: turn it into a standard FFT.}
Evaluating at the points $\omega_k$ is not yet a standard discrete Fourier
transform (DFT), whose points are the roots of unity $e^{2\pi ik/M}$. Two tricks
close the gap. First, since $\omega_k^{N/2}=e^{i\pi(4k+1)/2}=i$, the terms of
degree $j$ and $j+N/2$ can be grouped:
$p_j\omega_k^j+p_{j+N/2}\,\omega_k^{j+N/2}=(p_j+i\,p_{j+N/2})\,\omega_k^j$. This
\emph{folds} the $N$ real coefficients into $N/2$ complex numbers. Second, we
write $\omega_k^j=e^{i\pi j/N}\cdot e^{2\pi ijk/(N/2)}$. The first factor depends
only on $j$, and we multiply it into the folded vector (the \emph{twist}). The
second factor is a standard root of unity. Together,
\[
  p(\omega_k)=\sum_{j<N/2}\underbrace{\big(p_j+i\,p_{j+N/2}\big)\,e^{i\pi j/N}}_{z_j}
  \;e^{2\pi i jk/(N/2)},
\]
which is the standard DFT of size $N/2$ of the folded and twisted vector $z$
(Figure~\ref{fig:fft}). This is a special case of the weighted transforms of
\cite{crandall1994dwt}. To go back from values to coefficients, we do the same
steps in reverse order: inverse DFT, multiply by $e^{-i\pi j/N}$, and unfold the
real and imaginary parts.

% Folding and twisting: a negacyclic product of size N becomes a cyclic FFT of size N/2.
\begin{figure}[!htb]
\centering
\begin{tikzpicture}[font=\footnotesize, >=stealth,
    c/.style={draw, minimum width=0.5cm, minimum height=0.55cm, inner sep=0pt, fill=#1},
    c/.default=white,
    fstep/.style={->, line width=0.8pt},
    note/.style={font=\scriptsize, align=center, cGray!90!black}]
  \foreach \j in {0,...,3} \node[c=cBlue!12] (p\j) at (\j*0.5, 0) {$p_{\j}$};
  \foreach \j in {4,...,7} \node[c=cOrange!14] (p\j) at (\j*0.5, 0) {$p_{\j}$};
  \node[note, anchor=north] at (1.75,-0.4) {$N$ real coefficients};
  \begin{scope}[xshift=5.0cm]
    \foreach \j in {0,...,3} {
      \pgfmathtruncatemacro{\k}{\j+4}
      \node[c=cGreen!12, minimum width=1.2cm, font=\scriptsize] (z\j) at (\j*1.2, 0)
        {$p_{\j}+i\,p_{\k}$};
    }
    \node[note, anchor=north] at (1.8,-0.4) {$N/2$ complex values $z_j$};
  \end{scope}
  \draw[fstep] (3.95,0) -- node[above, note] {fold} (4.35,0);
  \node[draw, rounded corners=2pt, fill=cGray!8, align=center, inner sep=4pt]
    (tw) at (11.75,0) {twist $z_j\,e^{i\pi j/N}$\\[1pt]FFT of size $N/2$};
  \draw[fstep] (z3.east) -- (tw.west);
  \node[note, anchor=north] at (11.75,-0.75)
    {$=p(\omega_k)$, \ $\omega_k=e^{i\pi(4k+1)/N}$};
\end{tikzpicture}
\caption{Forward negacyclic transform, drawn for $N=8$. The two halves of the
coefficient vector are folded into $N/2$ complex numbers, which are multiplied
by the twist factors $e^{i\pi j/N}$. A standard FFT of size $N/2$ then gives
the values of $p$ at one root of $X^N+1$ from each conjugate pair. The inverse
transform does the same steps in reverse order.}
\label{fig:fft}
\end{figure}


\begin{example}[The complete computation for $N=4$]
We redo the previous example with the transform. To transform a polynomial
$p=[p_0,p_1,p_2,p_3]$ we fold it into $z=(p_0+ip_2,\;p_1+ip_3)$, multiply the
second entry by the twist $e^{i\pi/4}=0.707+0.707i$, and apply a DFT of size $2$,
which simply returns the sum and the difference of the two entries:
\begin{center}
\small
\begin{tabular}{@{}lccc@{}}
\toprule
 & folded $z$ & twisted & DFT: $(z_0+z_1,\;z_0-z_1)$\\
\midrule
$a=[1,1,0,0]$ & $(1,\;1)$ & $(1,\;0.707+0.707i)$ & $(1.707+0.707i,\;0.293-0.707i)$\\
$b=[0,0,0,1]$ & $(0,\;i)$ & $(0,\;-0.707+0.707i)$ & $(-0.707+0.707i,\;0.707-0.707i)$\\
\bottomrule
\end{tabular}
\end{center}
These are the values of $a$ and $b$ at $\omega_0$ and $\omega_1$ from the
previous table. Multiplying them one by one gives
$P=(-1.707+0.707i,\;-0.293-0.707i)$. To go back, the inverse DFT of size $2$
returns half the sum and half the difference, $(-1,\;-0.707+0.707i)$; undoing
the twist (multiplying the second entry by $e^{-i\pi/4}$) gives $(-1,\;i)$; and
unfolding, with real parts as $r_0,r_1$ and imaginary parts as $r_2,r_3$, gives
$r=[-1,0,0,1]$, the same result as the direct computation.
\end{example}

\paragraph{Step 4: the FFT itself.}
A DFT of size $M$ computed directly needs $M^2$ multiplications. The FFT of
\cite{cooley1965fft} is faster because it splits the problem: a DFT of size $M$
can be computed from two DFTs of size $M/2$, one on the entries with even index
and one on the entries with odd index, and those can be split again, until
size $1$. Figure~\ref{fig:butterfly} shows the case $M=4$. After putting the
inputs in a special order (bit-reversed), each round combines pairs of values
with a \emph{butterfly}, $(u,v)\mapsto(u+w\,v,\;u-w\,v)$, where $w$ is a root of
unity. For our size, $M=N/2=512$, this means $9$ rounds of $256$ butterflies,
that is, $2304$ butterflies instead of about $260{,}000$ multiplications. We
compute all the twist factors, roots of unity and the bit-reversed order once,
when the \code{NegFFT} object is created.

% Radix-2 FFT of size 4.
\begin{figure}[!htb]
\centering
% A DFT of size 4 computed with two rounds of butterflies (radix-2 FFT).
\begin{tikzpicture}[font=\footnotesize, >=Stealth,
    pt/.style={circle, fill=cNavy, inner sep=1.3pt},
    ln/.style={->, cNavy!70, line width=0.6pt, shorten >=2pt}]
  \def\dx{3.2}
  \foreach \lab [count=\r from 0] in {z_0,z_2,z_1,z_3}
    \node[pt, label=left:{$\lab$}] (a\r) at (0,-\r*0.85) {};
  \foreach \r in {0,...,3} \node[pt] (b\r) at (\dx,-\r*0.85) {};
  \foreach \lab [count=\r from 0] in {Z_0,Z_1,Z_2,Z_3}
    \node[pt, label=right:{$\lab$}] (c\r) at (2*\dx,-\r*0.85) {};
  % round 1: pairs (z0,z2) and (z1,z3)
  \foreach \u/\v in {0/1, 2/3} {
    \draw[ln] (a\u) -- (b\u);
    \draw[ln] (a\v) -- (b\u);
    \draw[ln] (a\u) -- (b\v);
    \draw[ln] (a\v) -- node[pos=0.78, below, font=\scriptsize, cOrange!90!black] {$-$} (b\v);
  }
  % round 2: pairs (b0,b2) and (b1,b3), with twiddles 1 and i
  \foreach \u/\v in {0/2, 1/3} {
    \draw[ln] (b\u) -- (c\u);
    \draw[ln] (b\v) -- (c\u);
    \draw[ln] (b\u) -- (c\v);
    \draw[ln] (b\v) -- node[pos=0.85, below, font=\scriptsize, cOrange!90!black] {$-$} (c\v);
  }
  \node[font=\scriptsize, cOrange!90!black, anchor=north] at (\dx,-3*0.85-0.08) {$\times i$};
  \node[font=\scriptsize, cGray, align=center] at (0.5*\dx,0.75) {round 1:\\sizes $1\to2$};
  \node[font=\scriptsize, cGray, align=center] at (1.5*\dx,0.75) {round 2:\\sizes $2\to4$};
  \node[anchor=west, align=left, font=\scriptsize] at (2*\dx+0.75,-1.27)
    {$Z_0=(z_0+z_2)+(z_1+z_3)$\\$Z_2=(z_0+z_2)-(z_1+z_3)$\\
     $Z_1=(z_0-z_2)+i\,(z_1-z_3)$\\$Z_3=(z_0-z_2)-i\,(z_1-z_3)$};
\end{tikzpicture}

\caption{A DFT of size $4$ computed by the FFT. The inputs are first put in
bit-reversed order ($z_0,z_2,z_1,z_3$). Each round combines pairs of values
with a \emph{butterfly}, $(u,v)\mapsto(u+w\,v,\;u-w\,v)$, where $w$ is a root of
unity ($1$, or $i$ for the last pair). A transform of size $M$ needs $\log_2M$
rounds of $M/2$ butterflies.}
\label{fig:butterfly}
\end{figure}


\paragraph{Is double precision enough?}
The FFT computes with \code{double} numbers, which have 53 bits of precision,
but the result must be exact as a 32-bit integer. In every product of the
scheme, one factor has small coefficients (digits with $|d|\le64$, or key bits)
and the other has torus coefficients, read as signed integers with
$|t|\le2^{31}$. Each coefficient of the product is a sum of $N=2^{10}$ such
terms, so its absolute value is at most $2^6\cdot2^{31}\cdot2^{10}=2^{47}$, which
is less than $2^{53}$. Rounding the result to the nearest integer therefore
gives the exact value. Our unit test confirms this: the FFT products were
identical, bit by bit, to the direct ones, and $\FftSpeedup\times$ faster
($\FftNaive$\,\textmu s against $\FftFast$\,\textmu s for $N=1024$).

\paragraph{A detail about complex numbers.}
The product of two \code{std::complex<double>} follows the C99 rules for
infinities and NaN values, which makes the compiler unable to vectorize it and
slows down the FFT loop. We use a hand-written function \code{cmul} with the
usual formula $(x+iy)(u+iv)=(xu-yv)+i(xv+yu)$.

\subsection{Signed Digits with One Addition: the Offset Trick}
\label{sec:offset}

\paragraph{The problem.}
Every external product decomposes $2N=2048$ coefficients into $2N\ell=6144$
digits, so the gadget decomposition must be fast. Ordinary digits in $[0,\Bg)$
would be easy: they are just groups of bits of $t$. But we want signed digits
in $[-\Bg/2,\Bg/2)$, because they are half as large ($\beta=\Bg/2$ instead of
$\Bg$), and the main noise term of Proposition~\ref{prop:extprod} grows with
$\beta^2$, so signed digits give four times less noise.

\paragraph{The slow way.}
To get signed digits from ordinary ones, we can work from the last digit to the
first, like in a manual subtraction: every digit that is $\Bg/2$ or larger
becomes negative by subtracting $\Bg$, and we add $1$ to the digit on its left.
Before that, we also have to round, using the part of $t$ that is thrown away.
Take base $10$ with two digits (digits in $[-5,5)$) and $t=0.9551$. Rounding to
two digits gives $0.96$, because the discarded part, $51$, is at least half.
The digits are $9$ and $6$. The last digit, $6$, is too large, so it becomes
$6-10=-4$ and the $9$ becomes $10$. Now $10$ is too large, so it becomes $0$ and
gives a $1$ to the left, which is a full turn of the circle and can be ignored.
The result is $(0,-4)$: indeed $0-0.04=-0.04\equiv0.96$. This works, but it needs
a loop with tests and carries, for each of the $6144$ digits.

\paragraph{The trick.}
The same result can be obtained with one addition. Add to $t$ a constant, the
\emph{offset}, that has $\Bg/2$ in every digit position plus half of the
position after the last digit; read the ordinary digits of the sum; and
subtract $\Bg/2$ from each of them. In base $10$ with two digits, the offset is
$0.555$. Table~\ref{tab:offset10} applies the trick to three values.

\begin{table}[!htb]
\centering
\caption{The offset trick in base $10$ with two digits: add $0.555$, keep the
part after the decimal point, take the two digits and subtract $5$ from each.}
\label{tab:offset10}
\small
\begin{tabular}{@{}lccccc@{}}
\toprule
$t$ & $t+0.555$ & digits & minus $5$ & value & error\\
\midrule
$0.7316$ & $1.2866\to0.2866$ & $2,\ 8$ & $-3,\ 3$ & $-0.27\equiv0.73$ & $0.0016$\\
$0.7366$ & $1.2916\to0.2916$ & $2,\ 9$ & $-3,\ 4$ & $-0.26\equiv0.74$ & $0.0034$\\
$0.9551$ & $1.5101\to0.5101$ & $5,\ 1$ & $0,\ -4$ & $-0.04\equiv0.96$ & $0.0049$\\
\bottomrule
\end{tabular}
\end{table}

The first row is Example~\ref{ex:decimal}. The second row shows the rounding:
the discarded part of $0.7366$ is $66$, so the value is rounded up to $0.74$.
The third row is the value we decomposed the slow way, and it gives the same
digits $(0,-4)$, including the chain of two carries. In every row the error is
at most $\varepsilon=0.005$.

\paragraph{Why it works.}
Adding $5$ to every digit and later subtracting $5$ from every digit does not
change the value. What changes is that the addition causes a carry exactly
when a digit was $5$ or more, and that is exactly when the slow way had to make
the digit negative and add $1$ to its left neighbor. The extra half digit at
the end does the rounding in the same way: it causes a carry into the last
digit exactly when the discarded part is at least half. The following lemma
states this for our binary case, where a digit is a group of
$\beta_{\mathrm{bit}}$ bits.

\begin{lemma}[Offset trick]
\label{lem:offset}
Let $q=2^{32}$, $\Bg=2^{\beta_{\mathrm{bit}}}$ with $\ell\beta_{\mathrm{bit}}<32$,
let $h_i=2^{32-i\beta_{\mathrm{bit}}}$ (the integer that stands for $\Bg^{-i}$),
and
\[
  \mathit{offset}=\sum_{i=1}^{\ell}\frac{\Bg}{2}\,h_i+\frac{h_\ell}{2}.
\]
For $t\in\Z_q$, let $t'=t+\mathit{offset}\bmod q$, let $u_i\in[0,\Bg)$ be the
$i$-th group of $\beta_{\mathrm{bit}}$ bits of $t'$ from the top, i.e.,
$u_i=\lfloor t'/h_i\rfloor\bmod\Bg$, and let $d_i=u_i-\Bg/2$. Then
$d_i\in[-\Bg/2,\Bg/2)$ and $\big|t-\sum_id_ih_i \bmod q\big|\le h_\ell/2$, so the
digits $(d_i)$ satisfy Eq.~\eqref{eq:gadget}.
\end{lemma}

\begin{proof}
We can write $t'=\sum_{i=1}^{\ell}u_ih_i+r$ with $0\le r<h_\ell$. Then, modulo $q$,
$t=t'-\mathit{offset}=\sum_i(u_i-\Bg/2)\,h_i+(r-h_\ell/2)=\sum_id_ih_i+(r-h_\ell/2)$,
and $|r-h_\ell/2|\le h_\ell/2$, which is $\varepsilon=1/(2\Bg^\ell)$ in torus
units. The range of $d_i$ follows from $u_i\in[0,\Bg)$.
\end{proof}

\paragraph{In binary.}
With $\Bg=2^7$ and $\ell=3$, a digit is a group of $7$ bits, and the offset has
$64$ in each of the three groups and $2^{10}$ in the discarded part.
Figure~\ref{fig:gadget} shows the computation for $t=0.7$: one 32-bit addition,
then each digit is obtained with a shift, a mask and a subtraction. The code
below is the complete decomposition of one coefficient; the offset is computed
only once.

% The offset trick for the signed gadget decomposition, on t = 0.7 with
% Bg = 2^7 and ell = 3.
\begin{figure}[!htb]
\centering
\begin{tikzpicture}[font=\footnotesize, >=stealth]
  \def\u{0.29}
  \def\h{0.5}
  % one 32-bit word split in fields of 7, 7, 7 and 11 bits (MSB on the left)
  \newcommand{\word}[6]{%
    \fill[#2] (0,#1) rectangle (21*\u,#1+\h);
    \fill[cGray!10] (21*\u,#1) rectangle (32*\u,#1+\h);
    \draw (0,#1) rectangle (32*\u,#1+\h);
    \foreach \b in {7,14,21} \draw (\b*\u,#1) -- (\b*\u,#1+\h);
    \node at (3.5*\u,#1+0.5*\h) {#3};
    \node at (10.5*\u,#1+0.5*\h) {#4};
    \node at (17.5*\u,#1+0.5*\h) {#5};
    \node at (26.5*\u,#1+0.5*\h) {#6};
  }
  \node[cGray!90!black] at (3.5*\u, 0.8) {$d_1$ field};
  \node[cGray!90!black] at (10.5*\u,0.8) {$d_2$ field};
  \node[cGray!90!black] at (17.5*\u,0.8) {$d_3$ field};
  \node[cGray!90!black] at (26.5*\u,0.8) {11 rounded bits};
  \word{0}{cBlue!10}{$89$}{$76$}{$102$}{$819$}
  \word{-0.7}{cGray!6}{$64$}{$64$}{$64$}{$2^{10}$}
  \draw[line width=0.6pt] (-0.1,-0.85) -- (32*\u+0.1,-0.85);
  \word{-2.1}{cGreen!12}{$26$}{$13$}{$38$}{$1843$}
  \word{-2.95}{cOrange!14}{$-38$}{$-51$}{$-26$}{---}
  % carries between fields
  \foreach \b in {21,14,7} {
    \draw[->, cRed, line width=0.7pt] (\b*\u+0.32,-1.5) to[bend right=50]
      node[above=-1pt, font=\scriptsize] {$+1$} (\b*\u-0.32,-1.5);
  }
  \draw[->, cRed, line width=0.7pt] (0.32,-1.5) to[bend right=50]
    node[above=-1pt, font=\scriptsize] {$+1$} (-0.32,-1.5);
  \node[cRed, font=\scriptsize, anchor=east] at (-0.38,-1.45) {lost (mod $1$)};
  % row labels
  \begin{scope}[every node/.style={anchor=east, font=\footnotesize}]
    \node at (-0.25, 0.25)  {$t=\texttt{0xB3333333}$};
    \node at (-0.25,-0.45)  {$+\ \mathit{offset}=\texttt{0x81020400}$};
    \node at (-0.25,-1.85)  {$t+\mathit{offset}=\texttt{0x34353733}$};
    \node at (-0.25,-2.7)   {digit $-\,B_g/2$};
  \end{scope}
  \node[anchor=north, align=center, font=\footnotesize] at (16*\u,-3.2)
    {$-\tfrac{38}{2^{7}}-\tfrac{51}{2^{14}}-\tfrac{26}{2^{21}}=-0.3000002\equiv0.6999998$,
     \quad error $1.9\cdot10^{-7}\le2^{-22}$};
\end{tikzpicture}
\caption{The offset trick for $t=0.7$, with $\Bg=2^7$ and $\ell=3$. The offset
has $64$ in each $7$-bit digit group and half of the last digit ($2^{10}$) in
the part that is thrown away. After one 32-bit addition, each digit is
obtained with a shift and a mask, and subtracting $64$ gives signed digits in
$[-64,64)$. The ordinary digits of $t$ are $(89,76,102)$; the carries of the
addition (red) make each digit that is $64$ or more negative and add $1$ to the
digit on its left.}
\label{fig:gadget}
\end{figure}


\begin{lstlisting}
uint32_t t = p[j] + g.offset;             // one addition: all carries
for (int i = 0; i < g.ell; ++i) {
    uint32_t digit = (t >> (32 - (i + 1) * g.Bgbit)) & g.mask; // [0,Bg)
    out[i][j] = int32_t(digit) - int32_t(g.halfBg);     // [-Bg/2,Bg/2)
}
\end{lstlisting}

The same idea of ``adding half before cutting'' is used in three other places:
to decode messages, in modulus switching (where $2^{31}$ is added to $2N\cdot t$,
computed in 64 bits, before dropping the 32 low bits), and in key switching,
where a half digit is added to each $a_i$ so that its $t$ digits are rounded
instead of truncated.

\subsection{External Product in the Frequency Domain}

The same rows of the bootstrapping key are used in every gate, so we store
them already transformed by the FFT: each TRGSW is kept as $2\ell$ rows of two
polynomials, each one as $N/2$ complex values. We also add up the products of
Eq.~\eqref{eq:extprod} while still in the frequency domain, where addition and
multiplication are done value by value (Figure~\ref{fig:extprod}). So an
external product costs $2\ell=6$ forward transforms (one per digit polynomial),
$4\ell=12$ value-by-value products, and only $2$ inverse transforms, about
$8\times630\approx5000$ FFTs of size $512$ per gate. The price is memory: a
complex \code{double} takes $16$ bytes for every two coefficients, twice as
much as the coefficients themselves, so the bootstrapping key takes
\BkMB\,MiB. Blind rotation also skips the steps where $\bar a_i=0$, since
$X^0-1=0$.

% Data flow of the external product with the bootstrapping key in the FFT domain.
\begin{figure}[!htb]
\centering
\begin{tikzpicture}[font=\footnotesize, >=stealth,
    box/.style={draw, rounded corners=2pt, fill=#1, minimum height=0.5cm,
                inner sep=2pt, align=center},
    box/.default=white,
    flow/.style={->, line width=0.6pt}]
  \node[box=cGray!8, minimum height=2.6cm, text width=1.3cm] (d) at (0,0)
    {TRLWE\\$d=(a,b)$};
  \node[box=cBlue!10, text width=1.6cm] (deca) at (2.45, 0.9) {decompose $a$};
  \node[box=cBlue!10, text width=1.6cm] (decb) at (2.45,-0.9) {decompose $b$};
  \draw[flow] (d.east |- deca) -- (deca);
  \draw[flow] (d.east |- decb) -- (decb);
  \foreach \i/\y in {1/1.5, 2/0.9, 3/0.3} {
    \node[box=cGreen!12] (fu\i) at (5.05,\y) {FFT$(u_{\i})$};
    \draw[flow] (deca.east) -- (fu\i.west);
  }
  \foreach \i/\y in {1/-0.3, 2/-0.9, 3/-1.5} {
    \node[box=cGreen!12] (fw\i) at (5.05,\y) {FFT$(w_{\i})$};
    \draw[flow] (decb.east) -- (fw\i.west);
  }
  \node[box=cOrange!14, minimum height=3.5cm, text width=2.15cm] (mac) at (8.35,0)
    {multiply and add\\value by value\\with the $2\ell$ rows\\of $\mathrm{BK}_i$\\(stored as FFT)\\[3pt]
     sums\\$\hat A,\ \hat B$};
  \foreach \n in {fu1,fu2,fu3,fw1,fw2,fw3} \draw[flow] (\n.east) -- (\n.east -| mac.west);
  \node[box=cGreen!12] (ia) at (11.15, 0.6) {IFFT$(\hat A)$};
  \node[box=cGreen!12] (ib) at (11.15,-0.6) {IFFT$(\hat B)$};
  \draw[flow] (mac.east |- ia) -- (ia);
  \draw[flow] (mac.east |- ib) -- (ib);
  \node[box=cGray!8, minimum height=2.6cm, text width=1.3cm] (r) at (13.4,0)
    {TRLWE\\$\mathrm{BK}_i\boxdot d$};
  \draw[flow] (ia.east) -- (ia.east -| r.west);
  \draw[flow] (ib.east) -- (ib.east -| r.west);
  \node[font=\scriptsize, cGray!90!black] at (5.05,-2.15) {$2\ell=6$ forward};
  \node[font=\scriptsize, cGray!90!black] at (8.35,-2.15) {$4\ell=12$ products};
  \node[font=\scriptsize, cGray!90!black] at (11.15,-2.15) {$2$ backward};
\end{tikzpicture}
\caption{One external product in our implementation. The digit polynomials
of both components are transformed once and multiplied by the rows of the
TRGSW ciphertext, which were transformed when the key was created. The
products are added up in the frequency domain, so each external product needs
only two inverse transforms. Blind rotation does one external product for each
nonzero $\bar a_i$, about $n=630$ per gate.}
\label{fig:extprod}
\end{figure}


\subsection{Key-Switching Table}

The key-switching key stores one ciphertext for every possible digit value, so
that key switching needs no multiplications: at most $Nt=8192$ subtractions of
vectors of $n+1$ torus elements, which take about \OurKS\,ms. The
$N\cdot t\cdot\Bks=32768$ ciphertexts, \KskMB\,MiB in total, are stored in a
single array, with entry $(i,j,v)$ at position $(i\,t+j)\,\Bks+v$, so that the
ciphertexts used for one coefficient are next to each other in memory. Entries
with $v=0$ are never used; they hold a trivial encryption of zero so that the
indexing stays simple.

%% ===========================================================================
\section{Experimental Evaluation}
\label{sec:eval}

The project proposal set four criteria for the comparison with the original
library: correctness (does every gate decrypt to the right bit?), performance
(how long does a gate take, and where is the time spent?), memory (how large
are the keys?) and noise (does the measured error match the theory?). This
section answers them in this order, after describing how we ran the
experiments.

\subsection{Setup}

We ran all measurements on a laptop with an \CPU{} processor, running Linux,
on a single thread fixed to one performance core, with the power profile set to
performance. Our code was compiled with \Compiler{} and the options
\code{-O3 -march=native}. The original library \cite{tfhelib} was compiled with
its own optimized configuration in two versions that differ only in the FFT:
\emph{spqlios-fma}, written in assembly with AVX2 and FMA instructions, and
\emph{nayuki-portable}, written in plain C. Both use the same 128-bit parameters
as ours (Table~\ref{tab:params}). The programs \code{apps/bench.cpp} and
\code{compare/bench\_original.cpp} do the same measurements with the two
libraries: the time to generate the keys, the average time of $500$ NAND gates,
the error of $500$ gates with random inputs (phase minus the expected
$\pm\frac18$), and the number of decryption errors. Each run creates new keys.
We did \OurRuns{} runs of our implementation and \OrigRuns{} of each version of
the original library, and we report the range of the results.

\subsection{Correctness}

All ten unit tests pass. The adder of Section~\ref{sec:adder} gave the right
sum in every run: \AdderEightText{} in \AdderEight\,s, \AdderSixteenText{} in
\AdderSixteen\,s, and \AdderTwentyFourText{} in \AdderTwentyFour\,s. The time
grows linearly with the number of bits, as expected from four blind rotations
per bit. Across the benchmark and noise experiments, \TotalGates{} gates with
random inputs gave no decryption error, and neither did the original library in
its \OrigGates{} gates. This agrees with the noise margin: the input of a NAND
is the sum of two bootstrapping outputs plus the modulus-switching error, with
a standard deviation of about $0.005$, while the margin is $\frac18=0.125$,
more than $20$ standard deviations away.

\subsection{Performance}

Table~\ref{tab:perf} and Figure~\ref{fig:time} summarize the comparison. Our
implementation computes a NAND gate in \OurNandMean\,ms, which is
$\SlowdownSpq\times$ slower than the original library with the assembly FFT and
$\SpeedupNay\times$ faster than the same library with the portable FFT. Since
the rest of the algorithm is the same in all three cases, this shows that the
speed of TFHE depends mostly on the quality of the FFT. Key generation is also
faster in our implementation. A likely reason is that the original library
creates and keeps the bootstrapping key both as coefficients and in the
frequency domain, while ours keeps only the second form.

\begin{table}[!htb]
\centering
\caption{Comparison with the original library (range over the runs, $500$ gates
per run, one thread). Memory is the peak memory used by the benchmark program.}
\label{tab:perf}
\small
\begin{tabular}{@{}lccc@{}}
\toprule
 & This work & \multicolumn{2}{c}{Original library}\\
\cmidrule(l){3-4}
 & & spqlios-fma & nayuki-portable\\
\midrule
Time per NAND gate (ms) & \OurNandRange & \SpqNandRange & \NayNandRange\\
Key generation (s) & \OurKeygenRange & \SpqKeygenRange & \NayKeygenRange\\
Output error variance ($10^{-5}$) & \OurVarRange & \SpqVarRange & \NayVarRange\\
Largest $|$mean$|$ of the output error & \OurMeanMax & \SpqMeanMax & \NayMeanMax\\
Decryption errors & 0 in \OurBenchGates & 0 in \OrigGatesEach & 0 in \OrigGatesEach\\
Peak memory (MiB) & \OurRSS & \SpqRSS & \NayRSS\\
\bottomrule
\end{tabular}
\end{table}

% Time per NAND gate and key generation time: this work vs. the reference library.
\begin{figure}[!htb]
\centering
\begin{tikzpicture}
\begin{groupplot}[
    group style={group size=2 by 1, horizontal sep=1.9cm},
    width=0.47\textwidth, height=5.4cm,
    ymin=0,
    xtick={0,1,2}, xticklabels={{this work},{spqlios-fma},{nayuki-portable}},
    xticklabel style={font=\footnotesize, align=center, text width=1.9cm},
    enlarge x limits=0.25,
    ymajorgrids, grid style={cGray!20},
    tick label style={font=\footnotesize}, label style={font=\small},
    title style={font=\small},
    nodes near coords, nodes near coords style={font=\scriptsize, anchor=south,
      yshift=1pt},
  ]
  \nextgroupplot[ybar, bar width=0.62cm, ylabel={time per NAND gate (ms)}, ymax=140,
                 title={(a) gate evaluation},
                 point meta=explicit symbolic]
  \addplot[fill=cBlue!70, draw=cBlue!90!black,
           nodes near coords={\pgfmathprintnumber[fixed,zerofill,precision=1]{\pgfplotspointmeta}}]
    table[x=x, y=nand, meta=nand,
          restrict expr to domain={\coordindex}{0:0}] {data/time.dat};
  \addplot[fill=cGray!45, draw=cGray!90!black,
           nodes near coords={\pgfmathprintnumber[fixed,zerofill,precision=1]{\pgfplotspointmeta}}]
    table[x=x, y=nand, meta=nand,
          restrict expr to domain={\coordindex}{1:2}] {data/time.dat};
  \nextgroupplot[ybar, bar width=0.62cm, ylabel={key generation (s)}, ymax=1.15,
                 title={(b) evaluation-key generation},
                 point meta=explicit symbolic]
  \addplot[fill=cBlue!70, draw=cBlue!90!black,
           nodes near coords={\pgfmathprintnumber[fixed,zerofill,precision=2]{\pgfplotspointmeta}}]
    table[x=x, y=keygen, meta=keygen,
          restrict expr to domain={\coordindex}{0:0}] {data/time.dat};
  \addplot[fill=cGray!45, draw=cGray!90!black,
           nodes near coords={\pgfmathprintnumber[fixed,zerofill,precision=2]{\pgfplotspointmeta}}]
    table[x=x, y=keygen, meta=keygen,
          restrict expr to domain={\coordindex}{1:2}] {data/time.dat};
\end{groupplot}
\end{tikzpicture}
\caption{Time per gate and key generation time, on one thread, with the same
parameters and on the same machine (average over the runs; the ranges are in
Table~\ref{tab:perf}). The original library is shown with its two FFT
versions; the rest of its code is the same in both.}
\label{fig:time}
\end{figure}


\paragraph{Where the time goes.}
To find the bottlenecks, we measured the time of each operation used by the
blind rotation on its own, and multiplied it by the number of times it runs
in one gate (Figure~\ref{fig:profile}). Blind rotation takes \BRShare\% of the
gate time and key switching the rest. The FFTs alone, forward and backward,
take \FftShare\% of the gate, and the value-by-value products another
\MacShare\%, so the work in the frequency domain is \FreqShare\% of the time.
The gadget decomposition, which with the offset trick is only a few integer
operations per coefficient, takes only \DecShare\%.

The measurements also showed a cost that does not appear in the description of
the algorithm: \emph{memory traffic}. Blind rotation reads a different TRGSW
ciphertext at each step, so every gate reads the whole \BkMB\,MiB bootstrapping
key, which is much larger than the processor caches. One external product takes
\ExtHot\,\textmu s when it reuses the same key ciphertext (which then stays in
the cache), but \ExtKey\,\textmu s when consecutive calls go through the key as
in the blind rotation. A single value-by-value product is $\MacRatio\times$
slower when its TRGSW row comes from main memory (\MacKey\,\textmu s instead of
\MacHot\,\textmu s), so Figure~\ref{fig:profile} counts the products with the
slower time. Over the $n$ steps, this difference is about \MemShare\% of the
gate time. The remaining \OtherShare\% (``other'') is the part that no single
operation explains: the temporary vectors created in every external product
(about \ExtAlloc\,\textmu s per call, measured as the difference between a full
external product that stays in the cache and the sum of its parts) and other
waits for memory. These results agree with the comparison above: the
$\SlowdownSpq\times$ difference to the original library comes from the work in
the frequency domain and from how it accesses memory, which the original
library implements in optimized assembly.

% Cost breakdown of one NAND gate, from the isolated timing of each primitive.
\begin{figure}[!htb]
\centering
\begin{tikzpicture}
\begin{axis}[
    xbar, bar width=0.42cm, width=0.6\textwidth, height=6.2cm,
    xmin=0, enlarge y limits=0.1, enlarge x limits={upper, value=0.22},
    xlabel={time per gate (ms)},
    ytick={0,...,6},
    yticklabels={{decomposition}, {rotation by $X^{\bar a_i}$, add/sub},
                 {key switching}, {other (temporaries, memory)},
                 {backward FFTs ($2n$)}, {value-by-value products ($4\ell n$)},
                 {forward FFTs ($2\ell n$)}},
    yticklabel style={font=\footnotesize},
    tick label style={font=\footnotesize}, label style={font=\small},
    xmajorgrids, grid style={cGray!20},
    point meta=explicit symbolic,
    nodes near coords={\pgfmathprintnumber[fixed,precision=1]{\pgfplotspointmeta}\,\%},
    nodes near coords style={font=\scriptsize, anchor=west},
  ]
  \addplot[fill=cBlue!70, draw=cBlue!90!black]
    table[x=ms, y=y, meta=pct] {data/profile.dat};
\end{axis}
\end{tikzpicture}
\caption{Where the time of a NAND gate goes. We measured each operation on its
own and multiplied its time by the number of times it runs per gate;
percentages are relative to the measured time of a gate. The value-by-value
products are measured with the TRGSW rows read from main memory, as in the
blind rotation. ``Other'' is the measured gate time minus all the listed
parts.}
\label{fig:profile}
\end{figure}


\subsection{Memory}

Most of the memory goes to the two parts of the evaluation key. The
bootstrapping key takes \BkMB\,MiB ($n\cdot2\ell\cdot2\cdot\frac N2$ complex numbers
of 16 bytes) and the key-switching key takes \KskMB\,MiB ($N\cdot t\cdot\Bks$
ciphertexts of $n+1$ torus elements of 4 bytes). An encrypted bit, on the other
hand, takes only $(n+1)\cdot4=2524$ bytes, which is still about $20{,}000$ times
more than the bit itself. The peak memory of our benchmark program is
\OurRSS\,MiB, less than the \SpqRSS\,MiB of the original library, which also
keeps the bootstrapping key as coefficients.

\subsection{Noise}
\label{sec:noise}

\paragraph{Bootstrapping resets the noise.}
First, a direct check of the property explained in Section~\ref{sec:bsnoise}.
We encrypted $+\frac18$ with a very large noise, more than $300$ times the noise
of a normal ciphertext, and bootstrapped each ciphertext once.
Figure~\ref{fig:noisereset} shows the result: after the bootstrapping, the
noise is the same as for any other input.

% Bootstrapping resets the noise (unit test test09).
\begin{figure}[!htb]
\centering
% Bootstrapping resets the noise: a deliberately noisy encryption of +1/8 and
% the output of its bootstrapping (unit test test09). Shared by report and slides.
\begin{tikzpicture}
\begin{axis}[width=11.5cm, height=4.6cm, xmin=0.075, xmax=0.175, ymin=0, ymax=150,
    axis lines=left, xlabel={phase}, ylabel={density},
    xtick={0.1,0.125,0.15}, xticklabels={$0.100$,$\frac18$,$0.150$},
    ytick=\empty, tick label style={font=\footnotesize}, label style={font=\small},
    legend style={font=\scriptsize, at={(0.99,0.97)}, anchor=north east, draw=none,
                  fill=white, fill opacity=0.85, text opacity=1},
    legend cell align=left, clip=false]
  \addplot[cOrange, line width=1.2pt, domain=0.075:0.175, samples=161]
    {exp(-(x-0.125)^2/(2*\BootSdInRaw^2))/(\BootSdInRaw*sqrt(2*pi))};
  \addlegendentry{input, $\sigma=\BootSdIn$}
  \addplot[cBlue, line width=1.2pt, fill=cBlue!15, domain=0.11:0.14, samples=161, visible on=<2->]
    {exp(-(x-0.125)^2/(2*\BootSdOutRaw^2))/(\BootSdOutRaw*sqrt(2*pi))};
  \addlegendentry{after bootstrapping, $\sigma=\BootSdOut$}
\end{axis}
\end{tikzpicture}

\caption{Bootstrapping resets the noise. Orange: the distribution of the phase
of encryptions of $+\frac18$ created with standard deviation \BootSdIn, more
than $300$ times the noise of a normal ciphertext. Blue: the phase after one
bootstrapping of each of them, with standard deviation \BootSdOut, the same as
for any other input.}
\label{fig:noisereset}
\end{figure}


\paragraph{Bounds and average case.}
The bounds of Eqs.~\eqref{eq:varbr} and~\eqref{eq:varks} are worst cases: they
assume that every digit has the largest possible size $\beta$ and that every
key-switching digit is nonzero. We get a better prediction by noting that the
digits are close to uniformly distributed. A uniform digit in $[-\beta,\beta)$
has $\mathbb{E}[d^2]\approx\beta^2/3$, so the average variance of the blind
rotation is $n\cdot2\ell N(\beta^2/3)\vartheta_{\BK}\approx\AvgBR$, a third of the
bound.

Key switching needs more care, because the table is created only once. For a
fixed key, each entry has a \emph{fixed} noise value. Let $X$ be the noise
added at position $(i,j)$, which is $e_{i,j,v}$ for a digit $v$ uniform in
$\{0,\dots,\Bks-1\}$, with $e_{i,j,0}=0$. Taking the average over the inputs
with the table fixed, and then over the random table, gives
\begin{equation}
  \mathbb{E}\big[\mathrm{Var}(X\mid\KSK)\big]
    =\Big(\frac{\Bks-1}{\Bks}\Big)^{2}\vartheta_{\KSK},
  \qquad
  \mathrm{Var}\big(\mathbb{E}[X\mid\KSK]\big)
    =\frac{\Bks-1}{\Bks^{2}}\,\vartheta_{\KSK}.
  \label{eq:ksbias}
\end{equation}
Adding the first term over the $tN$ positions predicts a variance of \AvgKS{}
for key switching within one key, so the total variance of a gate output, for
a fixed key, should be about \AvgTotal. The second term predicts that the
average error for a given key is not zero: each key has its own small
\emph{bias}, with standard deviation
$\sqrt{3tN/16}\,\sigma_{\mathrm{LWE}}\approx\BiasStd$ across keys.

\paragraph{Measurements.}
To test these predictions, we created \NoiseKeys{} independent keys and
bootstrapped \NoiseGates{} random NAND inputs with each of them
(Table~\ref{tab:noise} and Figure~\ref{fig:noise}). The measured variances are
far below the worst-case bounds and within a few percent of the average-case
predictions, both after blind rotation and after key switching. The standard
deviation of the average error of each key, \NoiseBiasStd, agrees with the
predicted \BiasStd{} (with only ten keys, this estimate has an uncertainty of
about $25\%$). The bias is harmless, since it is about a hundred times smaller
than the margin, but it explains why the average error of the original library
is smaller (Table~\ref{tab:perf}): when it creates the key-switching table, the
library draws all the noise values first and subtracts their average, which
removes most of the bias.

\begin{table}[!htb]
\centering
\caption{Variance of the output error of the gate bootstrapping: worst-case
bound, average-case prediction and measurement (\NoiseKeys{} keys $\times$
\NoiseGates{} gates; variance within each key, averaged over the keys).}
\label{tab:noise}
\small
\begin{tabular}{@{}lccc@{}}
\toprule
 & Bound & Average case & Measured\\
\midrule
After blind rotation and extraction & \BoundBR & \AvgBR & \NoiseVarBR\\
After key switching (within a key) & \BoundTotal & \AvgTotal & \NoiseVarKS\\
Average error of a key, std.\ deviation & --- & \BiasStd & \NoiseBiasStd\\
\bottomrule
\end{tabular}
\end{table}

% Output noise: histogram of the within-key error and per-key bias.
\begin{figure}[!htb]
\centering
\begin{tikzpicture}
\begin{groupplot}[
    group style={group size=2 by 1, horizontal sep=1.7cm},
    height=5.4cm,
    tick label style={font=\footnotesize}, label style={font=\small},
    title style={font=\small},
    grid style={cGray!20},
    scaled ticks=false,
  ]
  \nextgroupplot[width=0.56\textwidth, title={(a) error after key switching},
    xlabel={error $-$ mean of its key ($10^{-3}$)}, ylabel={density},
    xmin=-12, xmax=12, ymin=0, ymajorgrids,
    legend style={font=\scriptsize, at={(0.98,0.97)}, anchor=north east,
                  draw=cGray!50, fill=white, fill opacity=0.9, text opacity=1},
    legend cell align=left, yticklabels={}, ytick=\empty]
  \addplot[ybar interval=0.98, fill=cBlue!35, draw=cBlue!70!black]
    table[x expr=\thisrow{center}*1000-0.25, y=density] {data/noise_hist.dat};
  \addlegendentry{measured}
  \addplot[cOrange, line width=1.1pt, domain=-12:12, samples=161]
    {exp(-(x/1000)^2/(2*\AvgTotalRaw))/sqrt(2*pi*\AvgTotalRaw)};
  \addlegendentry{average-case prediction}
  \nextgroupplot[width=0.42\textwidth, title={(b) mean error of each key},
    xlabel={key}, ylabel={mean error ($10^{-3}$)},
    xmin=0.3, xmax=10.7, ymin=-3.2, ymax=3.2, xtick={1,...,10},
    ymajorgrids]
  \fill[cOrange!12] (axis cs:0.3,{-2000*\BiasStdRaw}) rectangle (axis cs:10.7,{2000*\BiasStdRaw});
  \fill[cOrange!25] (axis cs:0.3,{-1000*\BiasStdRaw}) rectangle (axis cs:10.7,{1000*\BiasStdRaw});
  \addplot[only marks, mark=*, mark size=2pt, cBlue]
    table[x=key, y expr=\thisrow{meanKS}*1000] {data/noise_keys.dat};
  \draw[cGray, densely dashed] (axis cs:0.3,0) -- (axis cs:10.7,0);
\end{groupplot}
\end{tikzpicture}
\caption{Output noise of the gate bootstrapping, for \NoiseKeys{} keys and
\NoiseGates{} gates per key. (a) Histogram of the error after key switching,
measured from the average error of its own key, compared with a Gaussian with
the predicted variance \AvgTotal. (b) The average error of each key; the shaded
bands show $\pm1$ and $\pm2$ predicted standard deviations of the bias of
Eq.~\eqref{eq:ksbias}. Both are very small compared with the margin
$\frac18=125\cdot10^{-3}$.}
\label{fig:noise}
\end{figure}


%% ===========================================================================
\section{Conclusion}
\label{sec:conclusion}

We implemented TFHE gate bootstrapping from scratch, in about \LibLines{} lines
of C++ without external libraries, and tested it layer by layer. With the
parameters of the original library, a gate takes \OurNandMean\,ms on one core.
This is $\SlowdownSpq$ times slower than the original library with its
assembly FFT and $\SpeedupNay$ times faster than its portable version, and our
implementation uses less memory. We saw no decryption error in \TotalGates{}
gates.

Building the scheme step by step made the role of each part clear. The
external product keeps the noise growth additive, which allows hundreds of
operations in one bootstrapping. The offset trick makes the signed
decomposition almost free. Most of the time is spent in the FFT and in reading
the bootstrapping key from memory, since the key is much larger than the
processor caches. Our noise measurements agree with an average-case analysis,
which also explained a small key-dependent bias in our key switching that the
original library removes by centering its noise.

There are several ways to continue this work, roughly from the easiest to the
hardest: centering the key-switching noise as the original library does;
removing the temporary vectors from the inner loops; storing the bootstrapping
key in one contiguous block and loading the next TRGSW ciphertext in advance;
using AVX2/FMA instructions in the FFT; running independent gates in parallel
threads; and using other test polynomials to compute more general functions
\cite{chillotti2021pbs}.

\bibliographystyle{sbc}
\bibliography{refs}

\end{document}
